\documentclass[conference,a4paper]{IEEEtran}

\usepackage{fontspec}
\setmainfont{Times New Roman}
\setsansfont{Arial}
\setmonofont{Courier New}[Scale=0.9]
\setlength{\emergencystretch}{2em}
\usepackage[english,russian]{babel}
\addto\captionsrussian{%
  \def\refname{Литература}%
  \def\figurename{Рис.}%
  \def\tablename{Таблица}%
  \def\abstractname{Аннотация}}
\usepackage{booktabs}
\usepackage{stfloats}
\usepackage[hidelinks]{hyperref}
\usepackage{pgfplots}
\pgfplotsset{compat=1.18}
\usetikzlibrary{arrows.meta,decorations.pathreplacing}
% Okabe-Ito palette, readable with colour vision deficiencies
\definecolor{okVermillion}{HTML}{D55E00}
\definecolor{okSky}{HTML}{56B4E9}
\definecolor{okOrange}{HTML}{E69F00}
\definecolor{okGreen}{HTML}{009E73}
\definecolor{okBlue}{HTML}{0072B2}
\DeclareUrlCommand\param{\urlstyle{tt}\def\UrlBreaks{\do\-\do\=}}

\title{Измерение накладных расходов\\механизмов доверия\\в слое обмена данными X-Road}
\author{\IEEEauthorblockN{Ruslan Nechshadimov}}

\newcommand{\newtext}[1]{{\color{blue}#1}}
\newcommand{\newred}[1]{{\color{red}#1}}
\newcommand{\newgreen}[1]{{\color{green!50!black}#1}}
\newcommand{\newviolet}[1]{{\color{violet}#1}}

\begin{document}
\maketitle
\thispagestyle{plain}
\pagestyle{plain}

\begin{abstract}
\newtext{X-Road защищает обмен данными между организациями с помощью подписей, проверки статуса
сертификатов, меток времени и журнала сообщений, и для оценки предлагаемых изменений этих механизмов,
например перехода на постквантовые подписи, нужна точка отсчёта: сколько они стоят сейчас. В работе
измеряются накладные расходы механизмов доверия X-Road 7.8.3 в расчёте на одно сообщение: на одном и
на трёх ноутбуках\newgreen{, на более быстром ноутбуке разработчика} и при добавленном времени
кругового обхода между серверами безопасности (Security Server); для этого используются изменения
конфигурации, поэтапные замеры времени и захват пакетов. При одном клиенте тракт X-Road
добавляет около 28\,мс к обмену размером 10\,КБ. На подпись и её проверку из них приходится лишь
несколько миллисекунд; б\'{о}льшая часть~--- это \newred{остаточная обработка} на сервере безопасности,
тогда как журналирование тела занимает почти половину обмена размером 1\,МБ, а синхронная простановка
меток времени добавляет около 90\,мс. При добавленном времени кругового обхода сильнее всего
сказывается открытие нового соединения для каждого сообщения, принятое по умолчанию: сообщение ожидает
около пяти круговых обходов при 10\,КБ и десяти при 1\,МБ против двух при повторно используемых
соединениях.} \newred{Оценка для постквантовых подписей сообщений, полученная из бенчмарка на
платформе X-Road и из сообщений, дополненных до соответствующего размера, показывает, что вычисления
меняются примерно на миллисекунду \newgreen{для схем на решётках и не более чем на 9\,мс для
составных схем с RSA или при SLH-DSA в удостоверяющем центре}, тогда как размер добавляет до
двух круговых обходов при ML-DSA и до пяти при SLH-DSA в удостоверяющем центре, причём только на новых
соединениях и главным образом за счёт подписей инфраструктуры. Таким образом, в рассмотренных
условиях задержка, добавляемая расстоянием, определяется новыми соединениями и окном TCP, а не
криптографией.}
\end{abstract}

\section{Введение}

Когда одна организация действует на основании данных, полученных от другой, она зависит от того, чего
сама не создавала и что не может проверить самостоятельно. Эту зависимость и называют доверием. В
исследованиях организаций доверие определяют как готовность одной стороны быть уязвимой к
действиям другой стороны, основанную на ожидании того, что та совершит действие, важное для
доверяющего, при отсутствии возможности наблюдать за ней или контролировать её~\cite{mayer1995}. В
вычислительной технике та же идея формулируется как убеждённость в способности сущности действовать
надёжно, безопасно и предсказуемо в заданном контексте~\cite{grandison2000}. В инженерии безопасности
\emph{доверенный} (trusted) компонент отличают от \emph{заслуживающего доверия} (trustworthy): первый~---
это компонент, отказ которого может нарушить политику безопасности, второй~--- компонент, который
не откажет~\cite{anderson2020}. Поэтому хорошо спроектированная система стремится требовать как можно
меньше доверия, а оставшееся сосредоточить в немногих компонентах, которые явно названы и
соответствующим образом защищены.

X-Road, слой обмена данными, построен на этом принципе. Организациям-участникам не нужно напрямую
доверять системам друг друга. Их доверие опирается на корень доверия~--- точку, которая принимается
доверенной аксиоматически~\cite{bakhtina2023sac}. В X-Road этот корень централизован: его образуют
удостоверяющие центры, ручающиеся за ключи участников~\cite{bakhtina2023sac}, вместе с
центральным сервером (Central Server), который публикует сведения о том, кто является участником и
каким центрам следует доверять, и службой меток времени~\cite{xroadarcss}. Каждое сообщение затем
несёт свидетельства, которые можно проверить относительно этого корня. Оно подписывается сервером
безопасности отправителя и проверяется сервером безопасности получателя, записывается в журналы
сообщений на обеих сторонах и снабжается меткой времени, а передаётся по каналу Transport Layer
Security (TLS) с взаимной аутентификацией, при этом статус сертификатов проверяется по протоколу
Online Certificate Status Protocol (OCSP)~\cite{xroadarcss,ansper2013}. Доверие к другой стороне
заменяется проверкой относительно корня доверия, и обмен впоследствии может служить
свидетельством того, кто, что и когда отправил. Обратная сторона централизованного корня в том, что он
становится единой точкой отказа~\cite{bakhtina2023sac}, поэтому и изучаются альтернативы ему.

Проверка не бесплатна. Подпись, её проверка, журналирование и установление канала требуют
вычислений и байтов: часть из них расходуется на каждое сообщение, часть~--- в фоне. X-Road измеряли и
раньше, но от начала до конца и ради других вопросов, а не ради стоимости этих механизмов: в
исследовании пороговых подписей сообщается о сквозных круговых обходах около 32\,мс при размещении
всех компонентов на одном ноутбуке, и они представлены как нижняя граница~\cite{bakhtina2026ijis}; в
исследовании пропускной способности измерены накладные расходы сервера безопасности перед одним
сервисом реестра, где преобладает время работы базы данных~\cite{taklai2020}; а в статье, в которой
была предложена пакетная простановка меток времени, дана аналитическая оценка, а измерение реализации
оставлено на будущее~\cite{ansper2013}. Данных о производительности, опубликованных Nordic Institute
for Interoperability Solutions (NIIS), который разрабатывает X-Road, найти не удалось. Ни один из
этих источников не разделяет стоимость обмена по механизмам.

Отсутствие такого разделения становится более существенным по двум причинам. Первая~--- объём:
пакетная простановка меток времени проектировалась в расчёте примерно на 240 миллионов запросов в
год~\cite{ansper2013}, тогда как эстонский экземпляр сейчас сообщает о 3.74 миллиарда~\cite{xteefactsheet},
так что любые расходы на одно сообщение умножаются на гораздо большее число. Вторая~--- ряд
предложений, которые сделали бы сами механизмы доверия дороже. Ожидается, что постквантовые подписи
увеличат сертификаты, подписи и рукопожатия~\cite{bakhtina2026pqc,snetkov2023}; например, подпись
Dilithium5 занимает 4595 байт против 256 байт у RSA-2048~\cite{snetkov2023}, \newtext{а в анализе
систем электронного правительства Эстонии утверждается, что подписи для X-Road должны быть
короткими из-за того, как там используется подпись~\cite{vakarjuk2024}}. Децентрализованная
инфраструктура открытых ключей (Public Key Infrastructure, PKI) на основе проверяемых учётных данных
добавляет к каждому запросу презентацию~\cite{bakhtina2023ares}, а пороговые подписи заменяют локальный ключ
распределённым протоколом~\cite{bakhtina2026ijis}. Эти работы описывают, что станет дороже, но без
измеренной точки отсчёта для нынешней системы не могут сказать, насколько.

{\color{blue}
Перечисленные предложения меняют не X-Road в целом, а отдельные механизмы, и большинство из них
меняет его криптографию. Постквантовые алгоритмы затрагивают подпись, её проверку и рукопожатие
TLS~\cite{bakhtina2026pqc}, а пороговые подписи меняют способ подписания
сообщений~\cite{bakhtina2026ijis}. Однако измерения защищённых протоколов показали, что стоимость
таких механизмов определяется не одной лишь криптографией. В анализе веб-серверов с TLS отдельные
операции TLS заменялись пустыми и измерялся выигрыш в пропускной способности сервера: операции RSA
оказались самым дорогим отдельным шагом, на них приходилось 13--58\,\% времени сервера, но даже при
их нулевой стоимости оставался большой разрыв с сервером без шифрования~\cite{coarfa2006}. Измерения
HTTPS показали, что стоимость TLS зависит главным образом от того, используются ли соединения
повторно~\cite{naylor2014}. Бенчмарки постквантовых схем в TLS показали, что на быстрых и надёжных
каналах время рукопожатия определяется вычислениями, тогда как при потере пакетов 3--5\,\% и более его
определяет размер сообщений~\cite{paquin2020}, и что цепочки сертификатов, которые больше не
помещаются в начальное окно перегрузки TCP, добавляют к рукопожатию круговые обходы, тогда как подпись
и проверка занимали менее 15\,мс у всех испытанных схем, кроме одной~\cite{sikeridis2020}.
\newtext{Эксперименты с увеличенными цепочками сертификатов подтвердили эту ступень на начальном
окне~\cite{westerbaan2021,chou2026}, а на соединениях, по которым передаются сотни килобайт, её доля в
общем времени снижается~\cite{kampanakis2024}.
\newtext{Бенчмарки на серверах, ноутбуках и периферийных устройствах показали, что серверы добавляют
менее 5\,\% задержки при постквантовом обмене ключами и подписях~\cite{abbasi2025}, а профилирование
TLS~1.3 с взаимной аутентификацией, ML-KEM и ML-DSA отнесло около 3\,\% процессорного времени на
криптографию, а остальное~--- на ожидание~\cite{rahul2026}. В микробенчмарках ML-KEM и ML-DSA
каждая операция занимала менее миллисекунды, тогда как в TLS~1.3 трафик рукопожатия от сервера вырос
с нескольких килобайт примерно до 8--15\,КБ, так что стоимость внедрения заключается в размере, а не в
вычислениях~\cite{demir2026}.}
Для подписей XML, формата сообщений X-Road, постквантовые и гибридные варианты были
реализованы с помощью Apache Santuario и BouncyCastle, и были измерены их подпись и проверка,
но цепочки сертификатов и сеть остались за рамками~\cite{mueller2024}.} X-Road проектировался с учётом
той же проблемы: пакетная простановка меток времени и кэширование ответов о статусе были введены
именно для того, чтобы запрос к доверенной службе на каждое сообщение не становился источником
задержки~\cite{ansper2013}. Вокруг своих криптографических операций X-Road, таким образом, принимает
архитектурные решения: он записывает каждое сообщение в журнал, проставляет метки времени на журнал
пакетами, а не на каждое сообщение, выносит проверку статуса сертификатов из тракта
запроса~\cite{ansper2013} и по умолчанию открывает новое соединение между серверами безопасности
для каждого сообщения~\cite{xroadsyspar}.

В совокупности эти работы делят стоимость механизма доверия на вычисления, определяемые его
алгоритмами, и часть, определяемую тем, как система их применяет: когда выполняются операции, что
сохраняется и как сообщения передаются по сети. Для X-Road это разделение не измерялось, а
перечисленные предложения изменили бы главным образом первую часть. Насколько существенно такое
изменение, зависит поэтому от того, насколько велика криптографическая часть стоимости по сравнению с
частью, определяемой архитектурой. Разделение может меняться и с условиями обмена. Более крупные
постквантовые подписи увеличивают сообщения, рост объёма означает больше одновременно обрабатываемых
сообщений, а самые свежие данные получены при размещении всех компонентов на одной машине, что их
авторы представляют как нижнюю границу~\cite{bakhtina2026ijis}, тогда как серверы безопасности двух
организаций обычно взаимодействуют по сети, где, согласно отчёту NIIS, постквантовые рукопожатия
должны добавить круговые обходы~\cite{bakhtina2026pqc}. Отсюда следуют исследовательские вопросы
(research questions, RQ):

\begin{description}
  \item[RQ1] Какая часть стоимости механизмов доверия X-Road в расчёте на одно сообщение приходится
  на криптографические операции, подпись и её проверку, а какая~--- на архитектурные решения
  вокруг них: журналирование сообщений, политику простановки меток времени и политику соединений?
  \item[RQ2] Как это разделение меняется с размером сообщения, с нагрузкой и с расстоянием между
  серверами безопасности?
  \item[RQ3] Сколько при таком разделении добавила бы на одно сообщение замена подписей RSA
  постквантовыми, в вычислениях и в круговых обходах, по сравнению с архитектурными расходами?
\end{description}

\newred{В этой работе под стоимостью механизмов доверия понимаются дополнительные вычисления, хранение
и обмен данными на один обмен, которые добавляют механизмы, устанавливающие, проверяющие и сохраняющие
свидетельства подлинности и целостности обмена X-Road; что в неё не входит, перечислено в
разделе~\ref{sec:bg}.}

Вклад работы~--- разделение стоимости доверия в X-Road, а не одна итоговая цифра бенчмарка.
Поскольку отключить все механизмы невозможно, криптографическая и архитектурная части
разделяются изменением конфигураций, поэтапной разбивкой, взятой из собственного операционного
мониторинга X-Road, и захватом пакетов, на стенде, который можно развернуть заново на одном или
нескольких узлах. Затем разделение исследуется в зависимости от размера сообщения, от нагрузки и от
управляемого времени кругового обхода, добавленного между серверами безопасности, где предсказание,
сделанное до измерения, было уточнено по захватам пакетов. Постквантовые схемы в X-Road пока
недоступны~\cite{bakhtina2026pqc,snetkov2023}, поэтому для RQ3 подпись и проверка с RSA и с
постквантовыми подписями измеряются вне X-Road \newtext{с той версией Java и той криптографической
библиотекой, которые использует X-Road (Java~21 и BouncyCastle~1.84),} и объединяются с измеренным
поведением TCP в оценку их стоимости на одно сообщение.\par}

\section{Предпосылки: за что платят на каждое сообщение}\label{sec:bg}

X-Road~--- централизованно управляемый распределённый слой обмена данными. Центральный сервер
распространяет глобальную конфигурацию~--- реестр участников, доверенные удостоверяющие центры (УЦ;
Certification Authority, CA) и доверенные службы меток времени,~--- а каждый участник эксплуатирует
сервер безопасности, который отвечает за управление ключами, подпись, защищённую передачу,
журналирование и контроль доступа~\cite{bakhtina2023ares,
xroadarcss}.

Для этой работы механизмы делятся по их положению относительно тракта запроса:

\begin{itemize}
  \item \textbf{Синхронные}, оплачиваемые на каждом сообщении: канал TLS с взаимной аутентификацией
  между серверами безопасности; подпись ключом подписи участника; проверка входящей
  подписи; запись запроса и ответа в журнал сообщений, причём тела по умолчанию журналируются
  целиком~\cite{xroadsyspar}. Сколько канал TLS стоит на одно сообщение, зависит от того,
  используются ли соединения повторно~\cite{naylor2014}. По умолчанию~--- нет: и повторное
  использование пула на стороне клиента (\param{pool-enable-connection-reuse}), и поддержка
  соединений из пула на стороне поставщика (\param{server-support-clients-pooled-connections})
  выключены, так что каждое сообщение открывает новое соединение TCP и выполняет рукопожатие
  TLS~\cite{xroadsyspar}.
  \item \textbf{Асинхронные}, амортизируемые или отложенные: проверка статуса сертификатов,
  использующая кэшированные ответы OCSP, которые обновляются с фиксированным интервалом (по умолчанию
  20 минут), а не запрос на каждое сообщение; простановка меток времени, выполняемая пакетами над
  хеш-цепочкой записей журнала (по умолчанию раз в минуту, при допустимом простое службы меток времени
  в четыре часа: параметр \param{acceptable-timestamp-failure-period} равен 14400 секундам); и
  периодическое обновление глобальной конфигурации~\cite{kivimaki2020,xroadsyspar,ansper2013}.
  Поскольку одна метка времени покрывает целый пакет, её стоимость на одно сообщение делится на размер
  пакета и снижается с ростом нагрузки~\cite{xroadarcss,ansper2013}.
\end{itemize}

\newtext{Подпись и проверка~--- это не по одной операции. В X-Road~7.8.3 каждый обмен
дважды подписывается ключом участника (запрос~--- потребителем, ответ~--- поставщиком) и
дважды проверяется. Каждое подписанное сообщение несёт, помимо значения подписи, сертификат подписи
и ответ OCSP для него, а запрос несёт ещё и ответ OCSP для сертификата аутентификации
потребителя. Получатель проверяет цепочку сертификатов сертификата подписи
на каждом сообщении, не кэшируя её; поставщик на каждом сообщении проверяет также цепочку сертификата
аутентификации потребителя, а потребитель цепочку поставщика~--- только при
открытии соединения. При одном уровне УЦ каждая цепочка стоит одной проверки подписи УЦ.
Подпись ответа OCSP и сертификат его ответчика проверяются один раз, результат
кэшируется на 60\,с, а УЦ и ответчик OCSP подписывают вне тракта запроса. Эти количества
установлены по исходному коду сервера безопасности\newtext{~\cite{xroadsrc}}.}
\newtext{На рис.~\ref{fig:path} показаны шаги одного обмена в том порядке, в котором их выполняют
серверы безопасности.}

\begin{figure*}[!tb]
\centering
\begin{tikzpicture}[font=\scriptsize, >={Latex[length=1.6mm]},
  box/.style={draw, rounded corners=1pt, minimum height=4mm, text width=5.0cm, align=center,
    inner sep=1.2pt, fill=white},
  crypto/.style={box, very thick, fill=okVermillion!45},
  conn/.style={box, dashed, fill=okOrange!45},
  cfg/.style={box, dashed, fill=okGreen!35},
  fix/.style={box, fill=okSky!35},
  msg/.style={->, thick},
  lbl/.style={above, fill=white, inner sep=1pt}]
% lanes
\foreach \x/\n in {0/Клиент, 3.4/Потребитель (ss1), 8.4/Поставщик (ss2), 12.3/Эхо-сервис} {
  \node[font=\footnotesize\bfseries] at (\x,0.5) {\n};
  \draw[densely dotted, black!50] (\x,0.3) -- (\x,-10.6);
}
% consumer: request
\draw[msg] (0,0) -- node[lbl]{запрос} (3.4,0);
\node[conn] at (3.4,-0.5) {новое соединение: TCP, возобновл. TLS};
\node[crypto] at (3.4,-1.0) {проверка цепочки аутент. поставщика (УЦ)};
\node[crypto] at (3.4,-1.5) {подпись запроса (участник)};
\node[fix] at (3.4,-2.0) {журнал: запись о запросе};
\node[cfg] at (3.4,-2.5) {журнал: тело запроса};
\draw[msg] (3.4,-3.1) -- node[lbl, pos=0.6]{подписанный запрос} (8.4,-3.1);
% provider: request
\node[crypto] at (8.4,-3.6) {проверка цепочки аутент. потребителя (УЦ)};
\node[crypto] at (8.4,-4.1) {проверка подписи (участник) и цепочки (УЦ)};
\node[fix] at (8.4,-4.6) {журнал: запись о запросе};
\node[cfg] at (8.4,-5.1) {журнал: тело запроса};
\draw[msg] (8.4,-5.65) -- node[lbl, pos=0.62]{пересылка} (12.3,-5.65);
\node[box, text width=1.6cm] at (12.3,-6.05) {сервис};
\draw[msg] (12.3,-6.45) -- node[lbl, pos=0.38]{ответ} (8.4,-6.45);
% provider: response
\node[crypto] at (8.4,-6.95) {подпись ответа (участник)};
\node[fix] at (8.4,-7.45) {журнал: запись об ответе};
\node[cfg] at (8.4,-7.95) {журнал: тело ответа};
\draw[msg] (8.4,-8.55) -- node[lbl, pos=0.4]{подписанный ответ} (3.4,-8.55);
% consumer: response
\node[crypto] at (3.4,-9.05) {проверка подписи (участник) и цепочки (УЦ)};
\node[fix] at (3.4,-9.55) {журнал: запись об ответе};
\node[cfg] at (3.4,-10.05) {журнал: тело ответа};
\draw[msg] (3.4,-10.6) -- node[lbl, pos=0.5]{ответ} (0,-10.6);
% stages of operational monitoring
\foreach \a/\b/\l in {0.05/-0.2/Отправка потребителем, -0.2/-3.1/До поставщика, -3.1/-5.65/Пересылка поставщиком,
  -5.65/-6.45/Эхо, -6.45/-8.55/Ответ поставщика, -8.55/-10.6/Завершение потребителем} {
  \draw[decorate, decoration={brace, amplitude=3pt}] (14.0,\a) -- (14.0,\b)
    node[midway, right=3pt, font=\scriptsize] {\l};
}
% legend and background
\node[crypto, text width=3.6cm] at (1.6,-11.25) {криптография};
\node[conn, text width=3.6cm] at (5.5,-11.25) {соединение (настраивается)};
\node[cfg, text width=3.6cm] at (9.4,-11.25) {журнал тел (настраивается)};
\node[fix, text width=3.6cm] at (13.3,-11.25) {\newred{остаточная обработка}};
\node[draw, densely dotted, rounded corners=1pt, text width=16.3cm, align=left, inner sep=3pt]
  at (8.1,-11.95) {В фоне, вне тракта запроса: ответы OCSP обновляются каждые 20 минут, а
  их проверка кэшируется на 60\,с; метки времени пакетами раз в минуту (\emph{synctsa}: по одной на
  запись журнала, в тракте); обновление глобальной конфигурации.};
\end{tikzpicture}
\caption{\newtext{Один обмен через X-Road~7.8.3 с конфигурацией по умолчанию, по исходному
коду сервера безопасности~\cite{xroadsrc}. Красный: криптографические операции; оранжевый и зелёный,
пунктир: шаги, которые меняет конфигурация (политика соединений и журналирование тел сообщений);
синий: \newred{остаточная обработка}, которая включает также разбор, формирование и пересылку сообщений. На повторно используемом соединении первые
два шага потребителя пропускаются. Участник: операция с ключом подписи участника; УЦ:
проверка подписи УЦ в цепочке сертификатов; аутент.: сертификат аутентификации. Фигурные скобки:
этапы, на которые операционный мониторинг X-Road делит обмен.}}
\label{fig:path}
\end{figure*}

Такое разделение~--- документированное проектное решение: кэширование ответов о статусе и отложенная
простановка меток времени были введены, чтобы операция PKI на каждое сообщение не становилась одновременно единой точкой отказа и
источником задержки~\cite{ansper2013}. У него есть и цена с точки зрения безопасности, поскольку
кэшированный ответ OCSP без nonce оставляет окно, в течение которого отозванный сертификат всё ещё
принимается~\cite{bakhtina2023ares}. Предмет этой работы~--- та сторона этого компромисса, которая касается производительности.

\newred{В измеряемую стоимость не входят доступность доверенных служб, выпуск
и отзыв сертификатов, обслуживание кэша OCSP в случаях, когда ответы истекают, отсутствуют
или ответчик недоступен, обновление глобальной конфигурации, полные рукопожатия TLS с
обменом сертификатами и управление ключами. Поэтому проверка статуса сертификатов измеряется в
установившемся режиме, когда кэш содержит действительные ответы.}

\section{Методология}

\subsection{Стенд}\label{sec:testbed}

На стенде X-Road~7.8.3 работает в контейнерах, собранных из официальных образов X-Road: один центральный сервер,
тестовый УЦ, который предоставляет также службы OCSP и меток времени, и четыре сервера безопасности~---
управляющий сервер, один потребитель (\emph{ss1}) и два поставщика (\emph{ss2}, \emph{ss3}). Измеряемый
тракт~--- \emph{генератор нагрузки $\rightarrow$ ss1 $\rightarrow$ ss2 $\rightarrow$ эхо-сервис}; второй
поставщик используется только для проверки того, что ограничивает потребителя под нагрузкой. Эхо-сервис отвечает на
REST-запрос телом запрошенного размера, заполненным случайными байтами, чтобы журнал сообщений не мог
его сжать, а собственное время обработки сервиса пренебрежимо мало. Генератор нагрузки~--- k6\newtext{~\cite{k6}}, он отправляет REST-запросы
\texttt{GET} в замкнутом цикле с заданным числом одновременных клиентов, так что размер сообщения
ниже~--- это размер тела ответа. Служба меток времени тестового УЦ была заменена
многопоточной реализацией, поскольку исходная обрабатывает по одному запросу за раз\newtext{, а
синхронная простановка меток времени отправляет запрос для каждой записи журнала сообщений, так что под нагрузкой
запросы выстраивались бы в очередь к службе и измерялась бы не политика}.

Каждый сервер безопасности ограничен 2 ядрами процессора и 4\,ГБ памяти, а центральный сервер~--- 2\,ГБ,
во всех окружениях, чтобы окружения различались тем, где работают контейнеры, а не тем, сколько
аппаратных ресурсов получает каждый из них. Весь стенд, включая регистрацию участников и сервисов,
разворачивается скриптами из одного файла композиции и может быть развёрнут заново с нуля.

\subsection{Окружения}

\paragraph{Три машины.} Три одинаковых ноутбука (Intel Core i5-8250U, 4 ядра и 8 потоков,
16\,ГБ памяти), соединённых выделенным проводным коммутатором. На первом работают центральный сервер, тестовый
УЦ, управляющий сервер и второй поставщик, на втором~--- потребитель и генератор
нагрузки, на третьем~--- поставщик и эхо-сервис. Контейнеры имеют собственные
адреса в сети коммутатора (Docker \texttt{ipvlan}), так что трафик между серверами
безопасности не транслируется узлом и покидает машину. Часы синхронизированы с помощью chrony\newtext{~\cite{chrony}} с точностью
до половины миллисекунды, что важно при сравнении меток времени с двух машин.
Схема показана на рис.~\ref{fig:testbed}.

\begin{figure}[t]
\centering
\begin{tikzpicture}[font=\footnotesize, >={Latex[length=1.6mm]},
  c/.style={draw, rounded corners=1pt, minimum height=5mm, align=center, inner sep=2pt, fill=white}]
\foreach \x/\n in {0/1, 2.85/2, 5.7/3} {
  \draw[rounded corners=2pt] (\x,0) rectangle ++(2.4,2.85);
  \node[anchor=north west, inner sep=2pt] at (\x,2.85) {\textbf{Ноутбук \n}};
  \draw (\x+1.2,0) -- (\x+1.2,-0.25);
}
\draw[fill=black!12] (0,-0.25) rectangle (8.1,-0.7);
\node at (4.05,-0.475) {Гигабитный коммутатор};
\node[c, minimum width=2cm] (cs) at (1.2,1.95) {Центр. сервер};
\node[c, minimum width=2cm] (ca) at (1.2,1.2) {УЦ, OCSP, TSA};
\node[c, minimum width=2cm] at (1.2,0.45) {ss0; ss3 + эхо};
\node[c, minimum width=1.5cm] (ss1) at (4.05,1.45) {ss1};
\node[c, minimum width=1.5cm] (k6) at (4.05,0.45) {k6};
\node[c, minimum width=1.5cm] (ss2) at (6.9,1.45) {ss2};
\node[c, minimum width=1.5cm] (echo) at (6.9,0.45) {эхо};
\draw[<->, very thick] (k6) -- (ss1);
\draw[<->, very thick] (ss1) -- node[above=1pt, fill=white, inner sep=1pt] {доб. RTT} (ss2);
\draw[<->, very thick] (ss2) -- (echo);
\draw[dashed] (cs.east) -- (2.6,1.95) -- (2.6,2.28) -- (6.9,2.28);
\draw[dashed, ->] (6.9,2.28) -- (ss2.north);
\draw[dashed, ->] (4.05,2.28) -- (ss1.north);
\end{tikzpicture}
\caption{Стенд на трёх машинах; тестовый УЦ предоставляет OCSP и службу меток времени (TSA). Жирная линия: измеряемый тракт; «доб. RTT»~--- добавленное время кругового обхода. Пунктирная линия: фоновый
трафик к центральным службам, который не задерживается. Весь трафик между серверами безопасности
проходит через коммутатор.}
\label{fig:testbed}
\end{figure}

\paragraph{Одна машина.} Тот же стек на одном из этих ноутбуков. При том же процессоре и
тех же ограничениях контейнеров отличие от окружения на трёх машинах состоит в сети между
серверами безопасности и в том, что один процессор делят все компоненты.

\paragraph{Ноутбук разработчика.} Тот же стек на более новом ноутбуке (Intel Core Ultra~7 265U) под
Windows~11, с контейнерами в WSL~2. Он показывает, как абсолютные значения зависят от процессора\newtext{.
Оценка для постквантовых подписей (RQ3, раздел~\ref{sec:pqmethod}) использует только этот ноутбук:
то, что в ней сравнивается, зависит от машины меньше, чем абсолютные задержки, поскольку каждая схема
сравнивается с RSA-2048 на том же процессоре, а круговые обходы подсчитываются по наклону зависимости от
добавленного времени кругового обхода, который определяется TCP, а не процессором}.

\paragraph{Эмулируемая сеть.} Во всех перечисленных окружениях серверы безопасности разделяет не более
одного коммутатора, и время кругового обхода составляет менее миллисекунды, тогда как серверы безопасности двух
организаций обычно взаимодействуют по более длинному и менее предсказуемому пути. По одним лишь локальным
измерениям нельзя сказать, сохранятся ли результаты там. Поэтому на трёх машинах между
потребителем и поставщиком добавляется управляемое время кругового обхода с помощью netem,
сетевого эмулятора ядра Linux~\cite{hemminger2005}, который добавляет выбранную задержку к пакетам
заданного потока. Это модель не реального развёртывания в глобальной сети, а только добавленного времени кругового обхода. Тот же подход используется для изучения TLS в разных сетевых условиях без
инфраструктуры масштаба Интернета~\cite{paquin2020}, а точность задержки, создаваемой netem,
исследована в~\cite{nussbaum2009,jurgelionis2011}. Как задержка применяется и проверяется,
описано в разделе~\ref{sec:delay}.

На всех машинах процессор работает в режиме питания «производительность», что проверяется перед каждой сеткой измерений и
записывается вместе с ней, поскольку режим питания влияет на задержку.

\subsection{Выделение стоимости механизма}\label{sec:iso}

Механизмы доверия из RQ1 нельзя отключать все по одному, поэтому каждый из них наблюдается
тем способом, который меняет как можно меньше остального (таблица~\ref{tab:mech}). \newtext{Архитектурные
решения из RQ1~--- это параметры конфигурации X-Road, поэтому их стоимость измеряется их изменением
при неизменных алгоритмах; подпись и проверку отключить нельзя, и они остаются в той
части тракта, которую конфигурация не меняет\newtext{, поэтому их стоимость отдельно измеряется вне
X-Road (раздел~\ref{sec:pqmethod})}.}

\begin{table}[t]
\centering\small
\caption{Как наблюдается стоимость каждого механизма доверия.}
\label{tab:mech}
\begin{tabular}{@{}p{0.40\columnwidth}p{0.54\columnwidth}@{}}
\toprule
Механизм & Через что наблюдается \\
\midrule
Весь тракт X-Road & \emph{direct} в сравнении с \emph{full} \\
Политика соединений между серверами & \emph{reuse} в сравнении с \emph{full} \\
Журналирование тел сообщений & \emph{nobody} в сравнении с \emph{full} \\
Подпись и проверка & этапы разбивки (раздел~\ref{sec:measured}) \\
Простановка меток времени & запросы к службе меток времени; \emph{synctsa} как другая политика \\
Проверка статуса сертификатов & запросы к ответчику OCSP \\
\bottomrule
\end{tabular}
\end{table}

Там, где настройка убирает или меняет механизм, конфигурация сравнивается с конфигурацией по умолчанию
(\emph{full}) и отличается от неё только этой настройкой:

\begin{itemize}
  \item \textbf{direct}~--- обычный HTTP от генератора нагрузки к эхо-сервису, без X-Road.
  Её отличие от \emph{full}~--- это весь тракт X-Road, а не какой-либо отдельный механизм.
  \item {\sloppy \textbf{reuse}~--- соединения между серверами безопасности остаются открытыми
  (\param{pool-enable-connection-reuse}, \param{server-support-clients-pooled-connections} и
  \param{server-connector-max-idle-time}). По умолчанию каждое сообщение открывает новое
  соединение~\cite{xroadsyspar}, так что разница~--- это эффект политики соединений между
  серверами безопасности: установление соединения TCP и сеанса TLS и начало каждой
  передачи на новом соединении TCP~\cite{naylor2014}. Рукопожатие TLS этим не выделяется.\par}
  \item \textbf{nobody}~--- тела сообщений не записываются в журнал сообщений
  (\param{message-body-logging=false}), так что разница~--- это стоимость журналирования тела.
  \item \textbf{synctsa}~--- на каждое сообщение метка времени проставляется синхронно
  (\param{timestamp-immediately=true}). Это другая политика, а не стоимость
  пакетной простановки меток времени по умолчанию, и именно так она приводится.
\end{itemize}

Подпись и проверку отключить нельзя, поэтому они наблюдаются только в тех этапах
разбивки, в которых происходят, вместе с остальной работой, выполняемой в этих этапах, а не как
разница между конфигурациями.

Два фоновых механизма нельзя перенести в тракт запроса, не изменив того, что
измеряется. Откладывание меток времени на длительный интервал нарушило бы документированное ограничение, связывающее
интервал простановки меток времени и свежесть OCSP~\cite{xroadsyspar}. Если сделать свежесть OCSP короче
интервала, с которым подписывающий компонент обновляет свои кэшированные ответы, каждый запрос \newtext{завершался бы} ошибкой,
поскольку сервер безопасности берёт статус сертификата только из этого кэша и не обращается к
ответчику во время обработки сообщения\newtext{~\cite{xroadsrc}}. Поэтому для этих двух механизмов стоимость наблюдается как число запросов,
отправленных службе меток времени и ответчику OCSP в ходе каждого прогона и пока на его сообщения
проставляются метки времени.

\subsection{Что измеряется}\label{sec:measured}

Для каждого запроса генератор нагрузки записывает задержку и то, пришёл ли корректный ответ
ожидаемого размера. Для каждого прогона записываются байты, переданные между потребителем и поставщиком
и подсчитанные \texttt{iptables} внутри сетевого пространства имён потребителя на порту X-Road,
прирост журнала сообщений и число запросов OCSP и запросов меток времени по журналу УЦ.

Кроме того, X-Road на каждом сервере безопасности записывает в базу данных операционного
мониторинга четыре метки времени на запрос: когда запрос поступил, когда он был переслан, когда пришёл ответ и
когда он был возвращён. Соединение записей потребителя и поставщика по идентификатору запроса
делит задержку на шесть последовательных этапов \newtext{(рис.~\ref{fig:path})}: потребитель до отправки, путь до поставщика
(этап,
в котором происходят подпись и журналирование на потребителе и установление соединения),
поставщик до пересылки (этап, в котором происходят проверка и
журналирование), эхо-сервис, возврат ответа поставщиком
и его завершение потребителем. В сумме этапы дают время между поступлением запроса на
потребитель и уходом ответа с него, куда не входит короткий путь между генератором нагрузки и
потребителем. Разбивка носит описательный характер:
это метки времени самой измеряемой системы с разрешением 1\,мс, и вместе с ней приводится доля
запросов, записи которых удалось сопоставить на обоих серверах. Этапы, измеренные в пределах
одного узла, надёжнее двух этапов, охватывающих разные узлы, где в результат входит остаточное
расхождение часов; объяснения, связанные с сетью, проверяются по захватам пакетов
(раздел~\ref{sec:delay}).

\subsection{Процедура}\label{sec:proc}

С одним клиентом используются два размера сообщений. Постквантовые подписи и сертификаты добавили бы
к каждому сообщению несколько килобайт~\cite{snetkov2023}, что существеннее всего для небольших сообщений, поэтому
ответ размером 10\,КБ представляет небольшое сообщение; ответ размером 1\,МБ представляет большое, тело которого журнал сообщений по
умолчанию сохраняет целиком~\cite{xroadsyspar}. Один клиент без другой нагрузки даёт базовую задержку без
очередей, а затем нагрузка увеличивается, пока у одного из компонентов не закончатся
ресурсы~\cite{taklai2020}: число клиентов превышает два ядра процессора, доступные каждому
серверу безопасности, и достигает 4, 16 и 32, чтобы выйти на потолок пропускной способности серверов безопасности; этот
диапазон охватывает 10--30 одновременных клиентов, использованных в~\cite{taklai2020}. Поэтому основная сетка измерений сочетает пять конфигураций с ответами размером 10\,КБ и 1\,МБ
при одном клиенте и 10\,КБ при 4, 16 и 32 клиентах, а также случай с двумя поставщиками при 4, 16 и 32 клиентах.

Сервер безопасности работает на виртуальной машине Java, первые запросы к которой после запуска выполняются медленнее,
пока код не скомпилирован, поэтому измеряется только установившийся режим~\cite{georges2007}. После каждого
изменения конфигурации прокси серверов безопасности перезапускается и прогревается отбрасываемым прогоном с четырьмя
клиентами, а перед каждым измеряемым прогоном выполняются прогревочные запросы. Объём прогрева задан по
пилотному прогону как точка, после которой медианная задержка по последовательным окнам запросов перестаёт
снижаться. В каждом прогоне измеряется 1000 запросов, как в~\cite{bakhtina2026ijis}, и 200 при 1\,МБ, где
каждое сообщение добавляет около мегабайта в журнал сообщений каждого сервера безопасности.

Следуя~\cite{georges2007}, каждая точка измеряется в нескольких независимых прогонах, а не один раз.
Сетка повторяет каждую точку пять раз; в каждом повторении сначала выполняется обычный HTTP, затем
конфигурации в случайном порядке, каждая со своими точками в случайном порядке, а контрольный прогон с конфигурацией
по умолчанию в начале каждого повторения и в конце показывает дрейф стенда за время
сетки. Выполняется несколько сеток, и их результаты объединяются. Во время измерений на машинах больше ничего не работает. Перед очисткой журнала сообщений
управляющий скрипт ждёт, пока каждое записанное сообщение не получит метку времени, чтобы отложенная работа не терялась
между прогонами.

Результаты приводятся как медиана по повторениям от медианы каждого прогона, с межквартильным
размахом между повторениями и числом успешных запросов.

\subsection{Добавленное время кругового обхода между серверами безопасности}\label{sec:delay}

Управляемое время кругового обхода добавляется между потребителем и поставщиком с помощью
netem~\cite{hemminger2005}, как это сделано для бенчмарков TLS в~\cite{paquin2020}. Половина
времени кругового обхода добавляется на каждой стороне, в сетевых пространствах имён двух серверов безопасности и только
для пакетов к другому серверу, так что задерживается весь трафик между ними, включая обмен OCSP
между серверами безопасности, а клиент, центральный сервер и УЦ остаются в локальной
сети. Задержка фиксированная, без джиттера, поскольку джиттер netem воспроизводит менее
точно~\cite{jurgelionis2011,nussbaum2009}. Результаты приводятся как эффект
управляемого добавленного времени кругового обхода между серверами безопасности, а не как измерение реального
развёртывания в глобальной сети. Контроль при нулевой задержке с теми же установленными правилами netem проверяет, что
сами правила не меняют результат.

Добавленное время кругового обхода (0, 5, 20 и 50\,мс) проверяется на самом измеряемом потоке: начиная со
второй сетки с задержкой время кругового обхода соединений потребителя с поставщиком считывается из
ядра во время каждого прогона, а калибровка с последовательными запросами считывает его для новых соединений.
Перед первой сеткой с задержкой захваты трафика последовательных запросов в локальной сети показывают,
открывает ли запрос новое соединение, с полным или возобновлённым рукопожатием TLS, или использует существующее,
и сколько раз одна сторона ждёт другую. По этим обменам делается предсказание того, как
задержка должна расти с временем кругового обхода, и оно сравнивается с измеренным наклоном; там, где они
расходятся, захваты при задержке используются для поиска ожиданий, которые предсказание упустило. Захваты ответа размером 1\,МБ
с повторным использованием соединений и без него проверяют, как передача большого ответа зависит от
состояния соединения TCP. Поскольку задержка применяется к исходящим пакетам, они ожидают внутри
отправляющего узла, где всё ещё могут учитываться в ограничениях стека TCP отправителя на данные в очереди
и тем самым влиять на рост его окна. Для перекрёстной проверки захваты при 1\,МБ и часть сетки
повторяются с задержкой, применяемой вместо этого к входящим пакетам на принимающей стороне.

Сетка с задержкой сочетает \emph{full} и \emph{reuse} с ответами размером 10\,КБ и 1\,МБ при одном клиенте
и 0, 5, 20 и 50\,мс, а расширенная сетка добавляет ответы размером 100\,КБ и 4 и 16 клиентов при 0, 20
и 50\,мс. Сетка перекрёстной проверки повторяет \emph{full} и \emph{reuse} с ответами размером 10\,КБ и 1\,МБ
при одном клиенте и 0 и 50\,мс.

\subsection{\newtext{Постквантовые подписи}}\label{sec:pqmethod}

{\color{blue}В X-Road~7.8.3 нет постквантовых подписей~\cite{bakhtina2026pqc,snetkov2023}, поэтому
измерить их на стенде напрямую нельзя. Новая схема подписи меняет в сообщении
две вещи: сколько времени занимают его подпись и проверка и сколько в нём байтов. Исследования
постквантового TLS разделяют эти два эффекта: вычисления определяют время на быстрых каналах, а
размер добавляет круговые обходы, как только данные перестают помещаться в окно TCP, которое на новом
соединении вначале мало и растёт с каждым круговым обходом~\cite{paquin2020,sikeridis2020}. Оценка для X-Road
следует тому же разделению и предполагает, что две части складываются. Подпись и проверка занимают
время процессора серверов безопасности, но не добавляют круговых обходов, поскольку подписи передаются
внутри сообщений, а не в отдельных обменах (раздел~\ref{sec:bg}); размер, в свою очередь,
меняет число круговых обходов, которое требуется для передачи сообщения на новом соединении, а его X-Road по умолчанию открывает
для каждого сообщения. Поэтому вычисления измеряются бенчмарком
схем подписи на той версии Java и с той криптографической библиотекой, которые использует X-Road, а
круговые обходы выводятся из размера сообщений и серий сегментов TCP (TCP flights), захваченных в сетках с
задержкой (раздел~\ref{sec:delay}). Часть, связанная с круговыми обходами, проверяется на стенде сообщениями, дополненными
до того же размера; предположение о том, что две части складываются, не проверяется. \newtext{Иными словами,
стоимость X-Road на одно сообщение раскладывается на слагаемые, и заменяются только те слагаемые, которые зависят от схемы
подписи: время подписи и проверки~--- временами новой схемы из бенчмарка,
а размер частей, относящихся к подписи,~--- их размером при новой схеме, влияние которого на круговые
обходы затем измеряется на X-Road сообщениями такого размера. Все остальные слагаемые, такие как журналирование,
канал TLS и пересылка, сохраняют значения, измеренные с RSA.} \newviolet{На рис.~\ref{fig:pqmodel}
показаны две части оценки и её проверка.}

\begin{figure*}[!tb]
\centering
\begin{tikzpicture}[font=\footnotesize,
  box/.style={draw, rounded corners=2pt, align=center, text width=2.5cm, minimum height=0.95cm, inner sep=2pt},
  comp/.style={box, fill=okVermillion!18},
  size/.style={box, fill=okOrange!22},
  check/.style={box, fill=okSky!25, dashed, text width=4.2cm},
  res/.style={box, fill=black!6, very thick, text width=3.7cm},
  arr/.style={-{Stealth[length=5pt]}, thick},
  row/.style={font=\footnotesize\bfseries, anchor=east, align=right}]
\node[row] at (-1.55,0) {Время};
\node[row] at (-1.55,-1.6) {Размер};
\node[comp] (c1) at (0,0) {Бенчмарк:\\подпись, проверка};
\node[comp] (c2) at (3.1,0) {$\times$ операций\\на обмен};
\node[comp] (c3) at (9.3,0) {Вычисления\\на сообщение};
\node[size] (s1) at (0,-1.6) {Сообщение сейчас\\(журнал сообщений)};
\node[size] (s2) at (3.1,-1.6) {Более крупные\\ключи и подписи};
\node[size] (s3) at (6.2,-1.6) {Серии\\сегментов TCP};
\node[size] (s4) at (9.3,-1.6) {Дополнительные\\круговые обходы};
\node[res] (e) at (13.1,-0.8) {Стоимость на сообщение:\\вычисления $+$\\круговые обходы $\times$ RTT};
\node[check] (k) at (6.2,-3.1) {Стенд: сообщения,\\дополненные до этого размера};
\draw[arr] (c1) -- (c2);
\draw[arr] (c2) -- (c3);
\draw[arr] (s1) -- (s2);
\draw[arr] (s2) -- (s3);
\draw[arr] (s3) -- (s4);
\draw[arr] (c3.east) -| ([xshift=-8pt]e.north);
\draw[arr] (s4.east) -| ([xshift=-8pt]e.south);
\draw[arr, dashed] (s2.south) |- (k.west);
\draw[arr, dashed] (k.east) -| node[pos=0, below right, inner sep=2pt, font=\scriptsize] {проверка} (s4.south);
\begin{scope}[shift={(10.6,-2.75)}, font=\scriptsize, sw/.style={draw, rounded corners=1pt, minimum width=0.45cm, minimum height=0.25cm, inner sep=0pt}]
\node[sw, fill=okVermillion!18] at (0,0) {}; \node[anchor=west] at (0.35,0) {вычисления (бенчмарк)};
\node[sw, fill=okOrange!22] at (0,-0.38) {}; \node[anchor=west] at (0.35,-0.38) {размер и круговые обходы (модель)};
\node[sw, fill=okSky!25, dashed] at (0,-0.76) {}; \node[anchor=west] at (0.35,-0.76) {измерено на X-Road};
\end{scope}
\end{tikzpicture}
\caption{\newviolet{Как оценивается стоимость постквантовой схемы подписи на одно сообщение. Время: времена подписи и проверки
каждой схемы из бенчмарка вне X-Road, умноженные на число этих операций в одном
обмене. Размер: части сообщения, относящиеся к подписи, с нынешними
ключами, увеличенные на ключи и подписи новой схемы, дают число серий сегментов TCP и тем самым
круговые обходы, добавляемые на новом соединении. Пунктир: круговые обходы проверяются на X-Road
сообщениями, дополненными до того же размера. То, что две части складываются, предполагается, но не проверяется.}}
\label{fig:pqmodel}
\end{figure*}

\newred{\paragraph{Рамки.} В отчёте NIIS названы пять областей, в которых X-Road предстоит миграция: TLS
между серверами безопасности и их клиентами, TLS между двумя серверами безопасности, подпись
сообщений, подписи глобальной конфигурации и подпись меток времени~\cite{bakhtina2026pqc}.
Оценка охватывает только подпись сообщений, включая подписи УЦ и ответчика OCSP,
которые несёт каждое сообщение (таблица~\ref{tab:pqscope}).}

\begin{table}[!htb]
\centering\footnotesize
\setlength{\tabcolsep}{3pt}
\caption{\newred{Рамки постквантовой оценки.}}
\label{tab:pqscope}
\color{red}
\begin{tabular}{@{}p{0.42\columnwidth}p{0.52\columnwidth}@{}}
\toprule
Часть миграции & В оценке \\
\midrule
Подписи участников на сообщениях & Да: бенчмарк, размер, дополненные сообщения \\
Подписи УЦ в сертификатах подписи и аутентификации & Да, как выше; один уровень УЦ \\
Ответы OCSP и сертификат ответчика & Да, как выше \\
TLS между серверами безопасности: обмен ключами, сертификаты рукопожатия & Нет: сеансы возобновляются; X-Road использует TLS~1.2 \\
TLS между серверами безопасности и клиентами & Нет \\
Подписи глобальной конфигурации & Нет: вне тракта запроса \\
Подписи меток времени & Нет: пакетами, вне тракта запроса \\
\bottomrule
\end{tabular}
\end{table}

\paragraph{Схемы и сценарии.} Рассматриваются схемы, названные в источниках по X-Road и по
постквантовому TLS~\cite{bakhtina2026pqc,snetkov2023,paquin2020,sikeridis2020}: ML-DSA на трёх её
уровнях безопасности~\cite{fips204}; Falcon-512 и Falcon-1024, которые NIST стандартизирует как FN-DSA;
составные подписи, сочетающие ML-DSA с RSA или ECDSA~\cite{rfc10073},~--- переход, рекомендуемый
отчётом NIIS~\cite{bakhtina2026pqc}; SLH-DSA~\cite{fips205}; и ECDSA, которую X-Road уже
поддерживает, как классическая альтернатива. Нынешняя конфигурация, ключи участников RSA-2048 под
УЦ с RSA-4096, служит базовой. Части сообщения, зависящие от схемы подписи, делятся на
две группы (раздел~\ref{sec:bg}): определяемые ключом участника~--- значение подписи и
открытый ключ в его сертификате, и определяемые ключами инфраструктуры~--- подпись
УЦ на сертификате участника, подпись ответа OCSP, а также ключ и подпись УЦ
в сертификате ответчика OCSP. Участники и инфраструктура могут мигрировать
по отдельности, поэтому сценарии различаются тем, какая из двух групп использует новую схему:
\begin{itemize}
  \item \emph{A}: меняются только ключи участников, УЦ и ответчик OCSP сохраняют RSA-4096;
  \item \emph{B}: участники, УЦ и ответчик OCSP используют новую схему;
  \item \emph{C}: участники используют ML-DSA, УЦ и ответчик OCSP~--- SLH-DSA;
  \item \emph{D}: меняются только УЦ и ответчик OCSP, участники сохраняют RSA-2048.
\end{itemize}
SLH-DSA рассматривается только для инфраструктуры (C и D). Её подписи занимают от 7.9 до 30\,КБ на
рассматриваемых уровнях~\cite{fips205} и передавались бы дважды в каждом обмене, а подпись у неё
медленная: для постквантовых подписей XML она была более чем в сто раз медленнее, чем у других
схем~\cite{mueller2024}, а при SLH-DSA в сертификате сервера рукопожатие TLS замедлялось
на порядки~\cite{delgado2026}. УЦ и ответчик OCSP подписывают редко и вне
тракта запроса (раздел~\ref{sec:bg}), так что там на каждое сообщение оплачивается только проверка их
подписей.

\paragraph{Вычисления.} \newtext{Поскольку эти схемы в X-Road недоступны, их подпись и
проверка измеряются бенчмарком вне его, но в той обстановке, в которой их выполнял бы сервер безопасности:
с BouncyCastle~1.84\newtext{~\cite{bouncycastle}} на Java~21, то есть с библиотекой и версией Java из X-Road~7.8.3; эта библиотека использовалась и для постквантовых подписей XML~\cite{mueller2024} и была быстрее
встроенной реализации Java~25 в более раннем исследовании на JMH~\cite{oleksiichuk2026}.
X-Road не подписывает само сообщение: сервер безопасности формирует
элемент \texttt{SignedInfo} подписи XML, который ссылается на сообщение только через его хеш (в
пакете~--- через результат хеш-цепочки), хеширует этот элемент и передаёт хеш подписывающему компоненту\newtext{~\cite{xroadsrc}}. Поэтому подписываемые данные малы и не растут с
сообщением, а хеширование самого сообщения одинаково для любой схемы; бенчмарк подписывает фиксированные
входные данные размером 1\,КБ, что соответствует порядку размера \texttt{SignedInfo}. Поскольку виртуальная машина Java компилирует код во время работы, подпись,
измеренная сразу после запуска, медленнее, чем в установившемся режиме работающего сервера
безопасности~\cite{georges2007}, поэтому в бенчмарке используется JMH\newtext{~\cite{jmh}}, который повторяет операцию после прогрева в
нескольких свежих виртуальных машинах: в 5 машинах по 10 односекундных итераций после 5 прогревочных итераций,
а для SLH-DSA-128s и 256s, у которых подпись медленная~\cite{mueller2024}, в 3 машинах по 10 пятисекундных
итераций. При одном клиенте операции обмена выполняются одна за другой в тракте запроса, и
сообщение ждёт каждую из них: потребитель подписывает запрос и, на новом соединении,
проверяет цепочку поставщика; поставщик проверяет подпись запроса и
цепочки сертификата подписи и сертификата аутентификации, прежде чем переслать запрос, и подписывает
ответ; потребитель проверяет ответ, прежде чем вернуть его (раздел~\ref{sec:bg}). Поэтому их времена
складываются, и стоимость обмена моделируется как сумма времён из бенчмарка по
этим операциям: две подписи и две проверки ключом участника и четыре проверки
подписи УЦ при \emph{full} и три при \emph{reuse}; проверка ответа OCSP
кэшируется и не учитывается. Эта сумма~--- также процессорное время, которое занимают подписи одного обмена.
Под нагрузкой сервер безопасности подписывает несколько ожидающих сообщений одной подписью над
хеш-цепочкой~\cite{xroadarcss}, так что стоимость подписи на одно сообщение снижается; оценка сделана для одного
клиента.} \newviolet{Поскольку подписывающий компонент получает только хеш, а ML-DSA и SLH-DSA в чистом
виде подписывают сами данные, X-Road потребовались бы их варианты с предварительным хешированием (pre-hash), HashML-DSA и
HashSLH-DSA~\cite{fips204,fips205}, и составные схемы с предварительным хешированием либо изменение этого интерфейса. Поэтому
чистые варианты и варианты с предварительным хешированием для ML-DSA, SLH-DSA-128f и трёх составных схем также
измеряются бок о бок в одном запуске с теми же настройками, чтобы увидеть, меняет ли выбор
варианта время; составная схема с предварительным хешированием получает хеш входных данных, как и подписывающий компонент
X-Road.}

\paragraph{Круговые обходы.} Дополнительные круговые обходы оцениваются в три шага:
\begin{enumerate}
  \item размер частей запроса и ответа, относящихся к подписи, с нынешними ключами
  (раздел~\ref{sec:bg}) измеряется в журнале сообщений стенда: подпись XML со
  значением подписи, сертификатом подписи и ответом OCSP внутри и ответ OCSP для
  сертификата аутентификации, который несёт запрос;
  \item по нему вычисляется их размер при новой схеме: новая схема увеличивает значение подписи,
  ключ участника и подпись УЦ в сертификате, а также подпись OCSP, ключ ответчика и
  подпись УЦ на сертификате ответчика, которые все закодированы в XML в base64,~--- каждый элемент на разность
  между закодированными размерами новой схемы и RSA;
  \item число серий сегментов TCP, которое требуется более крупным сообщениям на новом соединении, подсчитывается с помощью
  модели, в которой размеры серий (сколько сегментов несёт первая серия запроса и ответа
  и как растут следующие серии) взяты из захватов сеток с задержкой
  (раздел~\ref{sec:delay})\newgreen{, а там, где описанная ниже проверка расходится с моделью,~--- из
  захватов дополненных сообщений}.
\end{enumerate}
Модель проверяется на стенде дополнением тела запроса и ответа до размеров, которые она
предсказывает\newtext{. Постквантовые подписи увеличивают и запрос, и ответ, поскольку оба
несут подпись с сертификатом подписи и ответом OCSP (раздел~\ref{sec:bg}). Запросы
\texttt{GET} из остальных измерений не могут нести этот прирост: содержимое в запросе \texttt{GET}
не имеет определённого смысла в HTTP, и клиенту не следует его отправлять~\cite[разд.~9.3.1]{rfc9110}.
Поэтому для этой проверки генератор нагрузки отправляет запросы \texttt{POST}, тело которых имеет размер
прироста запроса, а эхо-сервис, возвращающий тело запрошенного размера
(раздел~\ref{sec:testbed}), отвечает ответом предсказанного размера}, при добавленном
времени кругового обхода 0 и 50\,мс, а для наибольших размеров~--- 0 и 20\,мс,
с тремя повторениями на точку и 150 запросами на прогон и с захватами трафика там, где модель
и стенд расходятся. \newviolet{Поскольку размеры серий, взятые из нескольких захватов, могут отличаться на сегмент,
что существенно только для сообщения вблизи границы серии, ещё одна сетка размещает точки,
запрос или ответ которых лежит в пределах примерно двух сегментов от границы модели, в том числе
точки сценария~D, при 0 и 50\,мс.} Задержка добавляется между двумя серверами безопасности так же, как на
трёх машинах (раздел~\ref{sec:delay}), и базовый обмен измеряется в той же сетке,
чтобы его наклон можно было сравнить с наклоном на трёх машинах. \newgreen{Дополнительные круговые
обходы точки~--- это её медианная задержка при добавленном времени кругового обхода минус задержка базового обмена
той же сетки, делённая на время кругового обхода, так что сдвиг всей машины между сетками
сокращается.}\par}

\section{Результаты}

\subsection{Стабильность стенда}

\newtext{В пилотном прогоне задержка не установилась в пределах первой тысячи запросов после
перезапуска серверов безопасности, первый прогон после перезапуска оставался более медленным даже после 4000
запросов, а повторные прогоны в одной и той же точке различались до 20\,\%; этим определены прогрев и
число повторений (раздел~\ref{sec:proc}). В энергосберегающем режиме питания ноутбуков, в
котором почти простаивающий процессор не повышает частоту для коротких всплесков работы, контрольный прогон занял
79.7 вместо 28.7\,мс; все приведённые ниже результаты получены в режиме производительности, а более ранние
прогоны в энергосберегающем режиме не используются.}

При пяти повторениях на сетку основная сетка была повторена 30 раз на трёх машинах (шесть сеток), 15
раз на одной машине (три сетки) и 20 раз на ноутбуке разработчика (четыре сетки). Контрольный прогон (\emph{full}, ответ 10\,КБ, один
клиент) имел медиану 28.7\,мс по 36 прогонам на трёх машинах (от 27.5 до 37.2\,мс, один выброс),
30.4\,мс по 18 прогонам на одной машине (от 28.9 до 31.2\,мс) и 26.7\,мс по 24 прогонам на ноутбуке
разработчика (от 24.3 до 27.9\,мс). Все запросы были успешными, кроме \emph{reuse} под нагрузкой, где единичные
запросы в некоторых прогонах завершались с HTTP~500: 7 из 90\,000 запросов при 4--32 клиентах на трёх
машинах, 3 из 45\,000 на одной машине и 2 из 60\,000 на ноутбуке разработчика. Эти отказы
согласуются с тем, что соединение из пула уже было закрыто другой стороной.

\subsection{\newtext{Разделение тракта при одном клиенте (RQ1)}}\label{sec:rq1}

В таблице~\ref{tab:onecl} приведена задержка при одном клиенте. Разница между \emph{direct} и
\emph{full}~--- это весь тракт X-Road: на трёх машинах 27.6\,мс для ответа 10\,КБ и 64\,мс для
ответа 1\,МБ. Журналирование тела сообщения при 10\,КБ не даёт разницы сверх разброса, но при
1\,МБ отказ от него экономит 34.9\,мс, 47\,\% задержки. Сохранение соединений между
серверами безопасности открытыми экономит 2.0\,мс при 10\,КБ и 4.3\,мс при 1\,МБ. Синхронная простановка меток времени добавляет
89--92\,мс. Межквартильный размах между повторениями составляет 0.7--4.8\,\% медианы на трёх машинах, до
9.4\,\% на ноутбуке разработчика и до 5.1\,\% на одной машине, кроме \emph{synctsa}
(18--31\,\%).

\begin{table}[!htb]
\centering\footnotesize
\setlength{\tabcolsep}{3.5pt}
\caption{Медианная задержка при одном клиенте (мс) по размеру ответа и окружению; разности с
\emph{full} в том же окружении; «Ноутбук»~--- ноутбук разработчика. Все запросы были успешными.}
\label{tab:onecl}
\begin{tabular}{@{}lrrrrrr@{}}
\toprule
& \multicolumn{2}{c}{3 машины} & \multicolumn{2}{c}{1 машина} & \multicolumn{2}{c}{Ноутбук} \\
\cmidrule(lr){2-3}\cmidrule(lr){4-5}\cmidrule(l){6-7}
& 10\,КБ & 1\,МБ & 10\,КБ & 1\,МБ & 10\,КБ & 1\,МБ \\
\midrule
\emph{direct} & 0.7 & 10.0 & 0.1 & 0.7 & 0.1 & 0.7 \\
\emph{full} & 28.3 & 74.0 & 29.8 & 74.7 & 24.8 & 49.7 \\
\emph{nobody} $-$ \emph{full} & $-0.9$ & $-34.9$ & $-0.3$ & $-32.0$ & $-0.9$ & $-17.0$ \\
\emph{reuse} $-$ \emph{full} & $-2.0$ & $-4.3$ & $-2.0$ & $-4.3$ & $-1.6$ & $-2.9$ \\
\emph{synctsa} $-$ \emph{full} & $+88.9$ & $+91.7$ & $+91.8$ & $+98.1$ & $+62.4$ & $+68.0$ \\
\midrule
повторения & \multicolumn{2}{c}{30} & \multicolumn{2}{c}{15} & \multicolumn{2}{c}{20} \\
\bottomrule
\end{tabular}
\end{table}

На трёх машинах обмен размером 10\,КБ передаёт 14.8\,КБ от потребителя к поставщику и 25.7\,КБ обратно;
при \emph{reuse}~--- 11.5 и 22.8\,КБ, так что политика соединений стоит 3.3 и 2.9\,КБ на сообщение в
двух направлениях. Журнал сообщений потребителя растёт на 33.7\,КБ на обмен с телами и на
19.3\,КБ без них, а при 1\,МБ~--- на 1.42\,МБ и 19.5\,КБ.

Два фоновых механизма в тракте запроса не проявились. Медианное число запросов к
ответчику OCSP за прогон было нулевым; примерно в 5\,\% прогонов (35 из 741 на трёх машинах) на
прогон пришлось фоновое обновление из 2--8 запросов. Конфигурация со свежестью OCSP
короче интервала обновления кэшированных ответов приводила к отказу каждого запроса, как описано в
разделе~\ref{sec:iso}; это наблюдалось в раннем прогоне на ноутбуке разработчика, но не зависит от
режима питания. Пакетная простановка меток времени отправляла службе меток времени в медиане ноль запросов
за время прогона (не более двух) и около двух на прогон, примерно из 1200 сообщений, пока управляющий скрипт ждал
меток времени; \emph{synctsa} отправляла 4802 на прогон, четыре на каждый обмен, по одному на
каждую запись, вносимую в журнал сообщений.

В таблице~\ref{tab:stages} приведена поэтапная разбивка на трёх машинах. Все запросы этих прогонов
были сопоставлены между потребителем и поставщиком. Этап до отправки запроса потребителем занимает
0.1--0.2\,мс, то есть потребитель фиксирует запрос как пересланный в момент начала отправки, а подпись и
журналирование запроса попадают в этап на пути к поставщику. Конфигурации меняют те
этапы, в которых работает их механизм: при 10\,КБ \emph{reuse} меняет только путь до поставщика ($-2.4$\,мс);
отказ от журналирования тела размером 1\,МБ сокращает ответ поставщика на 16.9\,мс, а этап завершения
у потребителя на 17.0\,мс, то есть журналирование тела стоит около 17\,мс на каждой стороне; а синхронная
простановка меток времени добавляет 20--25\,мс к каждому из четырёх этапов, в которых журналируется запись. Два этапа,
охватывающих разные машины, включают остаточное расхождение часов менее половины миллисекунды.

\begin{table*}[!tb]
\centering\small
\caption{Поэтапная разбивка по данным операционного мониторинга на трёх машинах, один клиент (среднее
по каждому прогону, медиана по 15 повторениям второй серии сеток, для которой сохранены записи обоих
серверов; мс; разрешение 1\,мс). Каждый столбец~--- собственная медиана, поэтому столбцы
не дают в сумме в точности итог. «До поставщика» и «Завершение
потребителем» охватывают две машины.}
\label{tab:stages}
\begin{tabular}{@{}lrrrrrrr@{}}
\toprule
& Отправка & До & Пересылка & Эхо & Ответ & Завершение & Итого \\
& потребителем & поставщика & поставщиком & & поставщика & потребителем & \\
\midrule
\emph{full}, 10\,КБ & 0.1 & 10.7 & 4.2 & 0.2 & 6.9 & 7.1 & 29.4 \\
\quad \emph{nobody} $-$ \emph{full} & 0.1 & 0.3 & 0.0 & 0.0 & $-0.8$ & $-0.7$ & $-1.2$ \\
\quad \emph{reuse} $-$ \emph{full} & 0.0 & $-2.4$ & 0.2 & 0.0 & 0.0 & 0.2 & $-2.1$ \\
\quad \emph{synctsa} $-$ \emph{full} & 0.1 & 24.3 & 20.6 & 0.2 & 19.8 & 25.1 & 90.5 \\
\midrule
\emph{full}, 1\,МБ & 0.2 & 12.4 & 5.4 & 0.3 & 33.0 & 24.5 & 76.5 \\
\quad \emph{nobody} $-$ \emph{full} & $-0.1$ & $-1.4$ & $-1.2$ & $-0.1$ & $-16.9$ & $-17.0$ & $-37.3$ \\
\quad \emph{reuse} $-$ \emph{full} & 0.0 & $-3.0$ & 0.0 & $-0.1$ & $-1.1$ & 0.0 & $-3.4$ \\
\quad \emph{synctsa} $-$ \emph{full} & 0.1 & 24.5 & 20.9 & 0.2 & 21.8 & 25.0 & 94.5 \\
\bottomrule
\end{tabular}
\end{table*}

{\color{blue}Таблица~\ref{tab:division} сводит эти результаты в разделение тракта X-Road. Расходы, которые
вытекают из архитектурных решений и могут быть изменены конфигурацией, малы для ответа
10\,КБ, 8--11\,\% тракта, и составляют 41--61\,\% тракта для ответа 1\,МБ, главным образом из-за журналирования
тела. То, что остаётся от тракта после них, составляет 22--27\,мс при 10\,КБ и вырастает на 0.1--10.3\,мс
при 1\,МБ, тогда как журналирование тела вырастает до 17--35\,мс. Этот остаток содержит подпись и
проверку на обоих серверах безопасности, но также журналирование записей о сообщениях без тел
(19.3\,КБ на обмен на потребителе), формирование и разбор сообщений и их пересылку, и
разделить их конфигурацией нельзя. Этапы, содержащие подпись или проверку, занимают
по 4.2--10.7\,мс при 10\,КБ, что ограничивает криптографические операции сверху, но не
измеряет их отдельно; отдельно их стоимость измеряется вне X-Road для RQ3
(раздел~\ref{sec:rq3}). Выбор между пакетной и синхронной простановкой меток времени~--- самое крупное
архитектурное решение в локальной сети: синхронная политика добавила бы 62--98\,мс, при 10\,КБ это
в \newgreen{2.5--3.2} раза больше всего тракта конфигурации по умолчанию.\par}

\begin{table}[!htb]
\centering\footnotesize
\setlength{\tabcolsep}{3.5pt}
\caption{Разделение тракта X-Road при одном клиенте (мс; по таблице~\ref{tab:onecl}). Весь тракт:
\emph{full} $-$ \emph{direct}; журналирование тела: \emph{full} $-$ \emph{nobody}; новое соединение:
\emph{full} $-$ \emph{reuse}; синхронная простановка меток времени, не входит в конфигурацию по умолчанию: \emph{synctsa} $-$
\emph{full}.}
\label{tab:division}
\color{blue}
\begin{tabular}{@{}lrrrrrr@{}}
\toprule
& \multicolumn{2}{c}{3 машины} & \multicolumn{2}{c}{1 машина} & \multicolumn{2}{c}{Ноутбук} \\
\cmidrule(lr){2-3}\cmidrule(lr){4-5}\cmidrule(l){6-7}
& 10\,КБ & 1\,МБ & 10\,КБ & 1\,МБ & 10\,КБ & 1\,МБ \\
\midrule
Весь тракт & 27.6 & 64.0 & 29.7 & 74.0 & 24.7 & 49.0 \\
\quad журнал тела & 0.9 & 34.9 & 0.3 & 32.0 & 0.9 & 17.0 \\
\quad новое соединение & 2.0 & 4.3 & 2.0 & 4.3 & 1.6 & 2.9 \\
\quad остаток & 24.7 & 24.8 & 27.4 & 37.7 & 22.2 & 29.1 \\
\midrule
Синхр.\ метки времени & $+88.9$ & $+91.7$ & $+91.8$ & $+98.1$ & $+62.4$ & $+68.0$ \\
\bottomrule
\end{tabular}
\end{table}


\subsection{Размер ответа и нагрузка (RQ2)}

При переходе от ответа 10\,КБ к ответу 1\,МБ задержка \emph{full} на трёх машинах растёт с 28.3 до
74.0\,мс, тогда как \emph{direct} растёт с 0.7 до 10.0\,мс. Без журналирования тела рост составляет от 27.4
до 39.1\,мс, так что б\'{о}льшая часть дополнительного времени при 1\,МБ~--- это запись тела в журнал сообщений.
Эффект политики соединений растёт с 2.0 до 4.3\,мс, а стоимость синхронной простановки меток времени
остаётся около 90\,мс. \newtext{Таким образом, с размером разделение смещается в сторону архитектуры:
остаток тракта вырастает не более чем на 10\,мс, тогда как журналирование тела вырастает до 17--35\,мс
(таблица~\ref{tab:division}).}

Под нагрузкой разброс между повторениями вырастает до 9--41\,\% пропускной способности для \emph{full},
\emph{nobody}, \emph{reuse} и случая с двумя поставщиками во всех окружениях (15--41\,\% на трёх
машинах), тогда как для \emph{synctsa} он остаётся на уровне 4--10\,\%, кроме четырёх клиентов на одной машине
(20\,\%).
Поэтому в таблице~\ref{tab:load} указана наибольшая пропускная способность, достигнутая каждой конфигурацией.
Различия между \emph{full}, \emph{nobody}, \emph{reuse} и двумя поставщиками лежат в пределах разброса
и не приписываются отдельным механизмам. Синхронная простановка меток времени вдвое снижает потолок на трёх
машинах (50 против 106 запросов в секунду) \newtext{и снижает его до 37\,\% на одной машине (28 против 75) и на ноутбуке разработчика (41 против 112)}. Второй поставщик также не поднимает потолок:
с двумя поставщиками этапы каждого поставщика становятся короче, но этап завершения у потребителя
становится длиннее, так что \newred{в этой конфигурации с двумя поставщиками} ограничением является потребитель.

\begin{table}[!htb]
\centering\footnotesize
\setlength{\tabcolsep}{4pt}
\caption{Наибольшая медианная пропускная способность, достигнутая при 4, 16 или 32 клиентах, ответ 10\,КБ (запросов в
секунду); в скобках число клиентов, при котором она достигнута, и межквартильный размах
пропускной способности между повторениями (\%).}
\label{tab:load}
\begin{tabular}{@{}lccc@{}}
\toprule
& 3 машины & 1 машина & Ноутбук \\
\midrule
\emph{full} & 106 (16; 32) & 75 (32; 29) & 112 (32; 9) \\
\emph{nobody} & 106 (32; 19) & 76 (16; 17) & 107 (32; 15) \\
\emph{reuse} & 108 (16; 21) & 79 (16; 20) & 118 (32; 15) \\
два поставщика & 111 (32; 24) & 82 (16; 36) & 109 (32; 24) \\
\emph{synctsa} & 50 (16; 10) & 28 (32; 5) & 41 (32; 5) \\
\bottomrule
\end{tabular}
\end{table}

\subsection{Окружение и добавленное время кругового обхода \newtext{(RQ2)}}\label{sec:rtt}

При одном клиенте три машины и одна машина дают одну и ту же картину (таблица~\ref{tab:onecl}):
задержка \emph{full} различается на 0.7--1.5\,мс, а разности между конфигурациями~--- не более чем на
6.4\,мс.
\emph{direct} показывает, что коммутатор добавляет 0.6\,мс при 10\,КБ и 9.3\,мс при 1\,МБ, но через X-Road
это компенсируется тем, что у поставщика собственный процессор: на трёх машинах этапы поставщика
короче (4.2 против 5.7\,мс и 6.9 против 9.0\,мс), а путь до поставщика на 0.6\,мс длиннее.
Под нагрузкой три машины достигают в 1.4--1.8 раза большей пропускной способности (таблица~\ref{tab:load}), поскольку
работа распределена по трём процессорам. При одном клиенте все абсолютные значения на ноутбуке разработчика ниже, а
расходы, зависящие от процессора, меньше (журналирование тела 1\,МБ~--- 17 вместо 35\,мс, синхронная простановка меток времени~--- 62
вместо 89\,мс), при этом разности сохраняют знак и порядок.

\paragraph{Добавленное время кругового обхода.} В трёх сетках с задержкой, в которых оно записывалось,
время кругового обхода соединений между серверами безопасности, сглаженное ядром (и
включающее отложенные подтверждения), составило 0.5, 6.1, 21.4 и 51.5\,мс для \emph{full} и 2.7, 8.1,
23.5 и 54.2\,мс для \emph{reuse} при добавленных 0, 5, 20 и 50\,мс; в калибровке с последовательными
запросами минимум для новых соединений составил 0.53, 5.54, 20.78 и 51.11\,мс. С установленными правилами netem и без задержки задержка ответа не изменилась при
10\,КБ (31.0 против 30.6\,мс для \emph{full}, 27.7 против 28.9\,мс для \emph{reuse}); при 1\,МБ она
была на 3--4\,мс выше (около 5\,\%, с перекрывающимися диапазонами), на один-два порядка меньше
эффекта добавленного времени кругового обхода (47\,мс для \emph{full} при 1\,МБ уже при 5\,мс).

Перед первой сеткой с добавленной задержкой захваты трафика в локальной сети показали, что каждое
сообщение \emph{full} открывает новое соединение TCP и возобновляет сеанс TLS~1.2 по сеансовому
билету (session ticket), без обмена сертификатами. Если считать рукопожатие TCP, рукопожатие TLS и
сам запрос, сообщение \emph{full} ожидает трёх круговых обходов между серверами безопасности,
а сообщение \emph{reuse}~--- одного. Поэтому рост задержки описывается её наклоном
относительно добавленного времени кругового обхода: наклон 3 означает, что задержка растёт на утроенное добавленное
время кругового обхода, то есть каждое сообщение ожидает трёх круговых обходов. Он вычисляется по медианам при 0 и
50\,мс. Предсказание давало наклон 3
для \emph{full} и 1 для \emph{reuse}. Измерение показано на рис.~\ref{fig:rtt} и в таблице~\ref{tab:rtt}: при 10\,КБ
наклон равен 5.1 для \emph{full} и 2.1 для \emph{reuse}, то есть предсказание занизило оба.
Захваты при задержке показали два ожидания, которые оно не учитывало: последняя часть запроса (842 байта
после примерно 10\,КБ) удерживается, пока поставщик не подтвердит предыдущие данные, что добавляет один круговой
обход в обеих конфигурациях, а на новом соединении ответ размером около 26\,КБ не помещается в
начальное окно TCP и требует второй серии сегментов, что добавляет к \emph{full} ещё один круговой обход\newgreen{:
ответ начинается с небольших записей TLS, начала сообщения MIME и разделителей
его фрагментов, отправляемых по одному пакету каждый, которые исчерпывают начальное окно, а тело следует только
после подтверждения}. С учётом
этого получается 5 и 2; калибровка с последовательными запросами на трёх машинах дала
наклоны 4.7 и 1.9. Короткое ожидание внутри рукопожатия TLS, между сообщением, которое переключает
потребителя на согласованные ключи, и его первым зашифрованным сообщением, проявилось при 20\,мс, но не при
50\,мс и вызывает колебание наклона примерно на 0.3--0.4. Поэтапная разбивка согласуется с этим: наклон
пути до поставщика равен 2.6 для \emph{full} и 0.5 для \emph{reuse}, а наклон этапа пересылки
поставщиком равен 1 в обеих конфигурациях.

\begin{figure}[t]
\centering
\begin{tikzpicture}
\begin{axis}[width=\columnwidth, height=6cm, xlabel={Добавленное время кругового обхода (мс)},
  ylabel={Медианная задержка (мс)}, xtick={0,5,20,50}, ymin=0, ymax=650, grid=major,
  legend style={font=\scriptsize, at={(0.5,-0.25)}, anchor=north, legend columns=2,
  /tikz/every even column/.append style={column sep=1em}}, legend cell align=left,
  tick label style={font=\small}]
\addplot[mark=*, thick, blue!70!black] coordinates {(0,28.5) (5,57.4) (20,131.1) (50,284.1)};
\addlegendentry{\emph{full}, 10\,КБ}
\addplot[mark=*, thick, orange!90!black] coordinates {(0,26.4) (5,37.5) (20,69.5) (50,131.0)};
\addlegendentry{\emph{reuse}, 10\,КБ}
\addplot[mark=square*, thick, densely dashed, blue!70!black] coordinates {(0,73.1) (5,120.0) (20,276.7) (50,608.4)};
\addlegendentry{\emph{full}, 1\,МБ}
\addplot[mark=square*, thick, densely dashed, orange!90!black] coordinates {(0,70.7) (5,81.5) (20,112.9) (50,172.0)};
\addlegendentry{\emph{reuse}, 1\,МБ}
\addplot[dotted, thick, blue!50, domain=0:50, samples=2] {28.5+3*x};
\addlegendentry{предсказание, \emph{full}, 10\,КБ}
\addplot[dotted, thick, orange!60, domain=0:50, samples=2] {26.4+1*x};
\addlegendentry{предсказание, \emph{reuse}, 10\,КБ}
\end{axis}
\end{tikzpicture}
\caption{Задержка при одном клиенте в зависимости от времени кругового обхода, добавленного между серверами безопасности,
три машины, 20 повторений; линии без
маркеров показывают предсказание, сделанное до измерения.}
\label{fig:rtt}
\end{figure}

\begin{table*}[!tb]
\centering
\caption{Точки рис.~\ref{fig:rtt}: медианная задержка по 20 повторениям (мс), в скобках
первый и третий квартили медиан прогонов.}
\label{tab:rtt}
\setlength{\tabcolsep}{4pt}
\begin{tabular}{@{}lcccc@{}}
\toprule
& 0\,мс & 5\,мс & 20\,мс & 50\,мс \\
\midrule
\emph{full}, 10\,КБ & 28.5 (27.9--28.7) & 57.4 (57.0--57.9) & 131.1 (130.9--131.4) & 284.1 (283.6--284.3) \\
\emph{reuse}, 10\,КБ & 26.4 (26.1--26.7) & 37.5 (37.0--37.6) & 69.5 (69.2--69.7) & 131.0 (130.7--131.1) \\
\emph{full}, 1\,МБ & 73.1 (71.5--76.1) & 120.0 (118.6--121.8) & 276.7 (275.6--281.1) & 608.4 (607.3--610.4) \\
\emph{reuse}, 1\,МБ & 70.7 (68.3--72.9) & 81.5 (80.0--82.8) & 112.9 (110.9--115.0) & 172.0 (170.3--173.7) \\
\bottomrule
\end{tabular}
\end{table*}

При 1\,МБ наклон равен 10.7 для \emph{full} и 2.0 для \emph{reuse}.
Захваты ответа размером 1\,МБ при 50\,мс показывают причину: на новом соединении тело уходит примерно
шестью сериями сегментов, по одной на круговой обход, каждая примерно вдвое больше предыдущей (28, 57, 113, 226 и
453\,КБ), по мере того как окно TCP растёт от начального размера; при \emph{reuse} окно сохраняется между
сообщениями и продолжает расти, так что ответ уходит меньшим числом серий, а начиная с пятого запроса~---
одной серией. Эффект
политики соединений, таким образом, вырастает с 2\,мс в локальной сети до 153\,мс (54\,\%) при 10\,КБ и
436\,мс (72\,\%) при 1\,МБ при добавленных 50\,мс.

Расширенная сетка подтверждает эту закономерность. При ответе 100\,КБ наклон равен 7.0 для \emph{full} и
2.1 для \emph{reuse}, то есть при новом соединении для каждого сообщения стоимость расстояния растёт с
размером ответа (наклоны 5.0, 7.0 и 9.7 для 10\,КБ, 100\,КБ и 1\,МБ в этой сетке), а
при повторно используемых соединениях~--- нет. Наклон при 1\,МБ здесь ниже, чем в первой сетке с задержкой
(9.7 против 10.7), поскольку медианы двух серий различаются на 8\,мс при 0\,мс и на 45\,мс при
50\,мс; при 10\,КБ две сетки согласуются (5.0 и 5.1). При добавленном времени кругового обхода
задержка с четырьмя клиентами та же, что с одним, а при 50\,мс пропускная способность ограничена круговыми обходами: 14 запросов в секунду для
\emph{full} против 30 для \emph{reuse}. С 16 клиентами при 0 и 20\,мс ограничением являются серверы
безопасности (103--127 запросов в секунду), и добавленное время кругового обхода растворяется в очереди (142
против 140\,мс для \emph{full}); при 50\,мс \emph{full} падает до 57 запросов в секунду, тогда как
\emph{reuse} сохраняет 104. Ошибки снова возникали только при \emph{reuse} под нагрузкой: 17 из 90\,000
запросов при 4 и 16 клиентах.

При задержке, применяемой к входящим, а не к исходящим пакетам, наклоны те же: 5.0 для
\emph{full} и 2.1 для \emph{reuse} при 10\,КБ и 9.6 и 2.0 при 1\,МБ (15 повторений, без
ошибок). Захваты ответа размером 1\,МБ при 50\,мс показывают те же серии на новом соединении (28, 57,
113, 226 и 453\,КБ), а при \emph{reuse} окно снова сохраняется между сообщениями. Следовательно, наклон
\emph{full} при 1\,МБ не зависит от того, где применяется задержка: он согласуется с
расширенной сеткой (9.7), тогда как первые сетки с задержкой дали 10.7, то есть от ночи к ночи он меняется примерно на один
круговой обход.

{\color{blue}С расстоянием разделение смещается ещё сильнее. Подписи передаются внутри сообщений, поэтому подпись и
проверка не добавляют собственных круговых обходов, а в поэтапной разбивке большинство круговых обходов приходится на
путь до поставщика и этап пересылки поставщиком, где устанавливается соединение и
передаётся запрос. Круговые обходы определяются политикой соединений и тем, как TCP передаёт
сообщение на новом соединении: при добавленных 50\,мс политика соединений стоит 153\,мс при 10\,КБ и
436\,мс при 1\,МБ против остатка около 25\,мс на тех же трёх машинах без добавленной задержки.

\subsection{\newtext{Постквантовые подписи (RQ3)}}\label{sec:rq3}

\newtext{В таблице~\ref{tab:pq} приведены бенчмарк и оценка для каждой схемы
(раздел~\ref{sec:pqmethod}). Одна подпись SLH-DSA занимала от 32\,мс до 1.1\,с. Проверка
RSA-4096, ключа тестового УЦ, занимала 85\,мкс. В журнале сообщений стенда
элемент \texttt{SignedInfo} занимал 813 байт в запросе и 825 байт в ответе, так что входные данные
бенчмарка близки к тому, что подписывает X-Road. На ноутбуке разработчика базовый обмен с
ответом 10\,КБ и \emph{full} имел наклон 5.4 относительно добавленного времени кругового обхода против 5.1 на
трёх машинах (таблица~\ref{tab:rtt}), и захваты на обоих показали одну и ту же последовательность
серий, так что круговые обходы, подсчитанные на ноутбуке, применимы и к трём машинам. Погрешность среднего (интервал 99.9\,\%) составила не
более 4\,\%, кроме подписи с ML-DSA и её составными схемами с ECDSA (от 11 до 35\,\%), поскольку подпись
с ML-DSA повторяет случайное число попыток.}

\newtext{С нынешними ключами подпись XML запроса занимала 7969 байт, подпись ответа~---
8002 байта, а ответ OCSP для сертификата аутентификации~--- 2052 байта. В захватах
сеток с задержкой новое соединение начиналось с окна в 10 сегментов,
первая серия \newgreen{базового} запроса после рукопожатия несла \newgreen{8} из них, начало ответа
состояло из небольших записей TLS, исчерпывавших начальное окно (раздел~\ref{sec:rtt}), а каждая
следующая серия удваивалась; \newgreen{первая версия} модели считает серии с этими размерами. При проверке модели круговых
обходов на стенде один запрос примерно из \newgreen{60\,600} завершился ошибкой. Наибольшему
запросу потребовалось на одну серию больше, чем предсказывала первая версия модели. Захват \newgreen{этого
дополненного запроса} показал, что
начиная с четвёртой серии запрос ограничен окном приёма поставщика (около
95\,КБ, то есть 36 и 41 сегмент вместо 80), а не окном перегрузки; модель была соответствующим образом
исправлена\newviolet{. Сетка вблизи границ серий показала ещё одно расхождение: запросу размером
41\,315 байт в канале (B, ML-DSA-65 + P-256) потребовались две серии, а не три. \newgreen{В
захватах дополненных запросов первая серия несла 9 полных сегментов. Базовый запрос помещается
в неё, но его последний, неполный сегмент уходил только после подтверждения, как и последняя часть
запроса в разделе~\ref{sec:rtt}, по причине, которую захваты не показывают; поэтому эта серия содержала только 8 полных
сегментов.} Теперь модель считает 9 сегментов в
первой серии и не менее двух серий на запрос}, и в таблице~\ref{tab:pqcheck}
исправленная модель сравнивается со стендом.}

\begin{table*}[!tb]
\centering\footnotesize
\caption{Схемы подписи на ноутбуке разработчика и их оценочная стоимость на одно сообщение относительно
базовой конфигурации (участники RSA-2048, УЦ RSA-4096). Вычисления (Выч.): изменение времени подписи и проверки
одного обмена при \emph{full} (базовое значение 1.7\,мс)\newgreen{, округлённое до 0.1\,мс; погрешность
среднего для подписи с ML-DSA и её составными схемами (интервал 99.9\,\%) составляет до $\pm$0.4\,мс на
обмен}. Размер: байты, добавляемые к запросу и к
ответу. Круговые обходы (RTT): дополнительные круговые обходы одного обмена при \emph{full}; при \emph{reuse} их нет
ни в одной строке. Подп. и Пров.: время одной подписи и одной проверки. Для SLH-DSA~--- сценарий C: участники ML-DSA-65, УЦ и ответчик OCSP~--- SLH-DSA.}
\label{tab:pq}
\color{blue}
\setlength{\tabcolsep}{3.5pt}
\begin{tabular}{@{}lrrrrrrrrrr@{}}
\toprule
& & & & & \multicolumn{3}{c}{A: ключи участников} & \multicolumn{3}{c}{B, C: участники, УЦ, OCSP} \\
\cmidrule(lr){6-8}\cmidrule(l){9-11}
Схема & Ключ (байт) & Подпись (байт) & Подп. (мкс) & Пров. (мкс) & Выч.\ (мс) & Размер (КБ) & RTT
& Выч.\ (мс) & Размер (КБ) & RTT \\
\midrule
RSA-2048 (базовая) & 294 & 256 & 668 & 26 & --- & --- & --- & --- & --- & --- \\
ECDSA P-256 & 91 & 70--72 & 54 & 52 & \newgreen{$-1.2$} & $-0.6$ / $-0.6$ & 0 & \newgreen{$-1.3$} & $-4.5$ / $-3.1$ & \newviolet{0} \\
ML-DSA-44 & 1334 & 2420 & 202 & 61 & \newgreen{$-0.9$} & 4.6 / 4.6 & 0 & \newgreen{$-1.0$} & 18.6 / 14.0 & 1 \\
ML-DSA-65 & 1974 & 3309 & 314 & 96 & \newgreen{$-0.6$} & 6.8 / 6.8 & 0 & \newgreen{$-0.5$} & 27.9 / 20.9 & 1 \\
ML-DSA-87 & 2614 & 4627 & 462 & 158 & \newgreen{$-0.1$} & 9.6 / 9.6 & 1 & \newgreen{$+0.1$} & 40.6 / 30.4 & 2 \\
Falcon-512 & 915 & 650--662 & 259 & 26 & \newgreen{$-0.8$} & 1.5 / 1.5 & 0 & \newgreen{$-1.1$} & \newgreen{3.3 / 2.7} & 0 \\
Falcon-1024 & 1811 & 1261--1277 & 519 & 52 & \newgreen{$-0.2$} & 3.7 / 3.7 & 0 & \newgreen{$-0.4$} & 11.6 / 8.8 & 1 \\
ML-DSA-44 + RSA-2048 & 1603 & 2676 & 881 & 90 & \newgreen{$+0.6$} & 5.4 / 5.4 & 0 & \newgreen{$+0.6$} & 21.6 / 16.2 & 1 \\
ML-DSA-44 + P-256 & 1398 & 2489--2492 & 359 & 114 & \newgreen{$-0.4$} & 4.8 / 4.8 & 0 & \newgreen{$-0.3$} & 19.4 / 14.6 & 1 \\
ML-DSA-65 + RSA-3072 & 2371 & 3693 & 2191 & 147 & \newgreen{$+3.3$} & 7.9 / 7.9 & 0 & \newgreen{$+3.5$} & 32.5 / 24.3 & 2 \\
ML-DSA-65 + P-256 & 2038 & 3378--3381 & 302 & 150 & \newgreen{$-0.5$} & 7.0 / 7.0 & 0 & \newgreen{$-0.2$} & 28.8 / 21.5 & \newviolet{1} \\
ML-DSA-87 + RSA-4096 & 3139 & 5139 & 4782 & 243 & \newgreen{$+8.7$} & 11.1 / 11.1 & 1 & \newgreen{$+9.3$} & 46.7 / 34.8 & 2 \\
ML-DSA-87 + P-384 & 2710 & 4729--4731 & 662 & 325 & \newgreen{$+0.6$} & 9.9 / 9.9 & 1 & \newgreen{$+1.5$} & 41.8 / 31.2 & 2 \\
SLH-DSA-128s & 50 & 7856 & 704\,547 & 619 & --- & --- & --- & \newgreen{$+1.6$} & 52.0 / 37.8 & 2 \\
SLH-DSA-128f & 50 & 17\,088 & 31\,957 & 1819 & --- & --- & --- & \newgreen{$+6.4$} & 110.3 / 77.6 & 4 \\
SLH-DSA-256s & 82 & 29\,792 & 1\,088\,102 & 1363 & --- & --- & --- & \newgreen{$+4.5$} & 190.6 / 132.5 & 5 \\
\bottomrule
\end{tabular}
\end{table*}

\begin{table}[!htb]
\centering\footnotesize
\setlength{\tabcolsep}{3.5pt}
\caption{Дополнительные круговые обходы при \emph{full}: модель в сравнении со стендом с дополненными сообщениями
(\newgreen{медианная задержка при добавленном времени кругового обхода минус задержка базового обмена той же сетки,
делённая на время кругового обхода}; одно значение на сетку).}
\label{tab:pqcheck}
\color{blue}
\begin{tabular}{@{}lccc@{}}
\toprule
Сценарий & Модель & 0 / 50\,мс & 0 / 20\,мс \\
\midrule
A ML-DSA-44 & 0 & 0.00 & \\
A ML-DSA-65 & 0 & $-0.33$ & \\
A ML-DSA-87 & 1 & 0.67 & \\
B ML-DSA-44 & 1 & 0.69 & \\
B ML-DSA-65 & 1 & 0.66 & \\
B ML-DSA-87 & 2 & 1.70, 1.67, 1.68 & 1.52 \\
B ML-DSA-87 + RSA-4096 & 2 & 1.69, 1.68 & 1.56 \\
C SLH-DSA-128s & 2 & 1.68, 1.69 & 1.56 \\
C SLH-DSA-128f & 4 & 3.74, \newviolet{3.70} & 3.85 \\
C SLH-DSA-256s & 5 & 4.65, 4.65 & 4.81 \\
\newviolet{B Falcon-1024} & \newviolet{1} & \newviolet{0.71} & \\
\newviolet{B ML-DSA-65 + P-256} & \newviolet{1} & \newviolet{0.71} & \\
\newviolet{D ML-DSA-44} & \newviolet{1} & \newviolet{0.72} & \\
\newviolet{D ML-DSA-44 + P-256} & \newviolet{1} & \newviolet{0.69} & \\
\newviolet{D ML-DSA-87} & \newviolet{2} & \newviolet{1.72} & \\
\bottomrule
\end{tabular}
\end{table}

\begin{figure*}[!tb]
\centering
\begin{tikzpicture}
\begin{axis}[name=fa, width=0.5\textwidth, height=5.2cm, xmode=log, log basis x=10,
  xmin=10, xmax=260, ymin=1.5, ymax=6.4, ytick={2,3,4,5,6},
  xtick={10,20,50,100,200}, xticklabels={10,20,50,100,200},
  xlabel={Байты в канале (КБ, логарифмическая шкала)}, ylabel={Серии сегментов TCP},
  title={(a) Серии по модели и точки сеток}, title style={font=\footnotesize},
  tick label style={font=\scriptsize}, label style={font=\footnotesize},
  legend style={font=\scriptsize, at={(0.03,0.97)}, anchor=north west}, legend cell align=left]
\addplot[const plot, very thick, okOrange] coordinates {(10,2) (41.992,3) (99.912,4) (152.04,5) (211.408,6) (260,6)};
\addlegendentry{запрос}
\addplot[const plot, very thick, okBlue, dashed] coordinates {(10,2) (28.96,3) (86.88,4) (202.72,5) (260,5)};
\addlegendentry{ответ}
\addplot[only marks, mark=*, mark size=1.6pt, okOrange, forget plot] coordinates {(12.6,2) (17.2,2) (19.4,2) (22.2,2)
  (31.1,2) (40.5,2) (53.2,3) (59.2,3) (64.5,3) (122.8,4) (203.2,5) (24.1,2) (41.3,2) (26.5,2) (27.1,2) (43.6,3)};
\addplot[only marks, mark=square*, mark size=1.4pt, okBlue, forget plot] coordinates {(20.7,2) (25.3,2) (27.5,2) (30.3,3)
  (34.7,3) (41.6,3) (51.1,3) (55.5,3) (58.5,3) (98.3,4) (153.2,4) (29.5,3) (42.2,3) (30.1,3) (30.5,3) (41.4,3)};
\node[font=\tiny, anchor=south east, inner sep=1pt] (lb) at (axis cs:36,2.45) {B ML-DSA-65 + P-256};
\draw[-{Stealth[length=3pt]}] (lb.south east) -- (axis cs:41.0,2.07);
\end{axis}
\begin{axis}[at={(fa.east)}, anchor=west, xshift=1.6cm, width=0.42\textwidth, height=5.2cm,
  xmin=-0.5, xmax=5.5, ymin=-0.6, ymax=5.3, xtick={0,1,2,3,4,5}, ytick={0,1,2,3,4,5},
  xlabel={Дополнительные круговые обходы, модель}, ylabel={Доп. круговые обходы, стенд},
  title={(b) Модель в сравнении со стендом (\emph{full})}, title style={font=\footnotesize},
  tick label style={font=\scriptsize}, label style={font=\footnotesize},
  legend style={font=\scriptsize, at={(0.03,0.97)}, anchor=north west}, legend cell align=left]
\addplot[fill=okSky!25, draw=none, forget plot] coordinates {(-0.5,-1) (5.5,5) (5.5,5.5) (-0.5,-0.5)} -- cycle;
\addplot[black!50, dashed, forget plot] coordinates {(-0.5,-0.5) (5.5,5.5)};
\addplot[only marks, mark=*, mark size=1.7pt, okOrange] coordinates {(0,0.00) (1,0.69) (0,-0.33) (1,0.66) (1,0.67)
  (2,1.70) (2,1.67) (2,1.69) (2,1.68) (4,3.74) (5,4.65) (2,1.68) (2,1.68) (2,1.69) (4,3.70) (5,4.65)};
\addlegendentry{RTT 50\,мс}
\addplot[only marks, mark=triangle*, mark size=2.2pt, okBlue] coordinates {(2,1.52) (2,1.56) (2,1.56) (4,3.85) (5,4.81)};
\addlegendentry{RTT 20\,мс}
\addplot[only marks, mark=diamond*, mark size=2.2pt, okGreen] coordinates {(1,0.71) (1,0.71) (1,0.72) (2,1.72) (1,0.69)};
\addlegendentry{вблизи границ серий, 50\,мс}
\end{axis}
\end{tikzpicture}
\caption{\newviolet{Модель серий сегментов TCP (раздел~\ref{sec:pqmethod}) и её проверка, новое
соединение (\emph{full}). (a) Серии, которые модель насчитывает для запроса (байты TLS после
рукопожатия) и для ответа (байты после его первой серии из небольших записей; эта серия
учитывается) заданного размера; точки~--- 16 размеров из сеток, каждый на своих двух линиях. Точка B
ML-DSA-65 + P-256 лежит на 0.7\,КБ ниже границы запроса; первая версия модели,
с первой серией в 8 сегментов, помещала её выше. (b) Дополнительные круговые обходы, насчитанные моделью,
в сравнении с измеренными: \newgreen{медианная задержка при добавленном времени кругового обхода минус задержка
базового обмена, делённая на время кругового обхода} (таблица~\ref{tab:pqcheck}); 26 точек.
Заливка: до половины кругового обхода ниже модели~--- та часть кругового обхода, которую базовый обмен уже
ожидает (раздел~\ref{sec:rtt}).}}
\label{fig:flights}
\end{figure*}

\begin{table}[!htb]
\centering\footnotesize
\setlength{\tabcolsep}{3.5pt}
\caption{\newtext{Сценарий D: меняются только УЦ и ответчик OCSP, участники сохраняют RSA-2048.
Столбцы как в таблице~\ref{tab:pq}.}}
\label{tab:pqD}
\color{blue}
\begin{tabular}{@{}lrrr@{}}
\toprule
УЦ и ответчик OCSP & Выч.\ (мс) & Размер (КБ) & RTT \\
\midrule
ECDSA P-256 & \newgreen{$-0.1$} & $-3.9$ / $-2.6$ & \newviolet{0} \\
ML-DSA-44 & \newgreen{$-0.1$} & 14.0 / 9.4 & 1 \\
ML-DSA-65 & \newgreen{0.0} & 21.1 / 14.1 & 1 \\
ML-DSA-87 & \newgreen{$+0.3$} & 31.0 / 20.7 & 2 \\
Falcon-512 & \newgreen{$-0.2$} & 1.8 / 1.2 & 0 \\
Falcon-1024 & \newgreen{$-0.1$} & 7.9 / 5.1 & 0 \\
ML-DSA-44 + RSA-2048 & \newgreen{0.0} & 16.2 / 10.9 & 1 \\
ML-DSA-44 + P-256 & \newgreen{$+0.1$} & 14.6 / 9.8 & 1 \\
ML-DSA-65 + RSA-3072 & \newgreen{$+0.3$} & 24.5 / 16.3 & 1 \\
ML-DSA-65 + P-256 & \newgreen{$+0.3$} & 21.7 / 14.5 & 1 \\
ML-DSA-87 + RSA-4096 & \newgreen{$+0.6$} & 35.5 / 23.7 & 2 \\
ML-DSA-87 + P-384 & \newgreen{$+1.0$} & 31.9 / 21.3 & 2 \\
SLH-DSA-128s & \newgreen{$+2.1$} & 45.2 / 31.0 & 2 \\
SLH-DSA-128f & \newgreen{$+6.9$} & 103.5 / 70.8 & 4 \\
SLH-DSA-256s & \newgreen{$+5.1$} & 183.8 / 125.7 & 5 \\
\bottomrule
\end{tabular}
\end{table}

\newtext{Вычисления меняются мало:} с ECDSA,
ML-DSA-44 и ML-DSA-65, Falcon и составными схемами ML-DSA-44 и ML-DSA-65 с ECDSA
подписи одного обмена занимают на 0.2--1.3\,мс меньше, чем с RSA-2048, поскольку подпись RSA~---
самая дорогая операция базовой конфигурации, а с ML-DSA-87~--- примерно столько же. \newgreen{Составная схема
ML-DSA-44 с RSA-2048 добавляет около 0.6\,мс.} Составная схема
ML-DSA-87 с P-384 добавляет от 0.6 до \newgreen{1.5}\,мс, составные схемы с RSA-3072 и RSA-4096, в подписи которых
преобладает половина RSA, добавляют от 3.3 до 9.3\,мс, а SLH-DSA в роли УЦ добавляет от 1.6 до 6.4\,мс за счёт
проверки подписи УЦ на каждом сообщении. Вся подпись и проверка
базовой конфигурации, 1.7\,мс,~--- это 7\,\% от 24.7\,мс, которые тракт X-Road занимает на том же ноутбуке
(таблица~\ref{tab:division}). \newtext{Эти времена относятся к ноутбуку разработчика. На трёх
машинах расходы, зависящие от процессора, были в 1.4--2.1 раза выше (синхронная простановка меток времени 89
против 62\,мс, журналирование тела 1\,МБ 35 против 17\,мс, таблица~\ref{tab:onecl}). При масштабировании на эти
коэффициенты подписи базовой конфигурации заняли бы там от 2.4 до 3.5\,мс, а составные схемы с
RSA-3072 и RSA-4096 и SLH-DSA в УЦ добавили бы до 13--19\,мс вместо не более 9.3\,мс,
что сопоставимо с остатком тракта в 24.7\,мс на тех же машинах; для ML-DSA и Falcon
изменение остаётся в пределах примерно 1\,мс.}

\newviolet{Варианты с предварительным хешированием занимают столько же времени, сколько чистые. В запуске, где те и другие
измерялись бок о бок, проверка различалась не более чем на 7\,мкс (0.4\,\% для SLH-DSA-128f, от 1 до 5\,\% для
ML-DSA и составных схем), подпись с ML-DSA и составными схемами~--- не более чем на 83\,мкс, в пределах
погрешности среднего для этой подписи, которая повторяет случайное число попыток, а подпись
SLH-DSA-128f, которую выполняют только УЦ и ответчик OCSP,~--- на 0.9\,\%. Для двух
подписей и двух проверок ключом участника в одном обмене это менее 0.2\,мс, так что
оценки таблицы~\ref{tab:pq}, сделанные с чистыми вариантами, справедливы и для вариантов с предварительным хешированием, которые
потребовались бы интерфейсу подписывающего компонента X-Road.}

Размер имеет значение только на новом соединении. Когда меняются только ключи участников (A), сообщения вырастают на
1.5--11\,КБ и не требуют дополнительного кругового обхода, кроме ML-DSA-87 \newgreen{и двух её составных схем}, где ответ уже не
помещается в свою вторую серию. Когда меняются также УЦ и ответчик OCSP (B), растут и сертификаты,
и ответы OCSP в каждом сообщении, и сообщения вырастают на 3--47\,КБ: один дополнительный
круговой обход с ML-DSA-44, ML-DSA-65, Falcon-1024\newviolet{,} составными схемами ML-DSA-44 \newviolet{и
ML-DSA-65 + P-256}, два~--- с ML-DSA-87\newviolet{, ML-DSA-65 + RSA-3072} и составными схемами
ML-DSA-87\newviolet{; с ECDSA, которая уменьшает сообщения, число серий остаётся прежним}. Falcon-512~--- единственная постквантовая схема, которой они не требуются \newtext{ни в одном} сценарии. С SLH-DSA
для УЦ сообщения вырастают на 52--191\,КБ и требуют от двух до пяти дополнительных круговых обходов.
\newtext{Когда меняются только УЦ и ответчик OCSP (D, таблица~\ref{tab:pqD}), сообщения вырастают
почти так же, как в B: с ML-DSA~--- на 14--31\,КБ в запросе, около трёх четвертей
прироста в B, поскольку подпись УЦ и ответчика OCSP встречается в каждом сообщении трижды:
в сертификате подписи, в ответе OCSP и в сертификате ответчика,
тогда как ключ участника встречается только в его подписи и сертификате. Круговые обходы те же,
что в B, кроме Falcon-1024 и \newviolet{ML-DSA-65 + RSA-3072}, которым требуется на один меньше, а
вычисления меняются менее чем на 1\,мс, кроме SLH-DSA (от 2.1 до 6.9\,мс).} При
\emph{reuse} окно уже открыто, и ни одной строке не требуется дополнительный круговой обход.

\newtext{Все} \newviolet{26} \newtext{точек таблицы~\ref{tab:pqcheck} согласуются с исправленной моделью} \newviolet{(рис.~\ref{fig:flights})} при округлении
до целых круговых обходов, при 50 и при 20\,мс, а ступени кратны добавленному времени кругового обхода
(SLH-DSA-256s против 128f: 45\,мс при 50\,мс и 19\,мс при 20\,мс). Измеренные значения лежат на 0--0.5
кругового обхода ниже модели, поскольку сам базовый обмен уже ожидает часть кругового обхода
до подтверждения своей первой серии, прежде чем отправить последние байты запроса
(раздел~\ref{sec:rtt}). При \emph{reuse} \newgreen{ни одна точка не отстояла от базового обмена более чем на 0.41 кругового
обхода} ни в одном сценарии,
а две точки, повторённые из первой ночи, различались на 3--5\,мс. \newgreen{Задержка
ноутбука разработчика без добавленной задержки не была стабильной от сетки к сетке: контрольный прогон занимал
от 29 до 45\,мс в зависимости от сетки и менялся до 19.5\,\% в пределах одной сетки. Поэтому значения при 0\,мс
для проверки не используются; сдвиг всего ноутбука одинаково влияет на базовый обмен и на точку
при добавленном времени кругового обхода и сокращается в их разности.}

\newred{Из этих результатов времена бенчмарка и задержка дополненных сообщений на X-Road
измерены; вычисления на обмен, размер постквантовых сообщений и число серий
вычислены по ним с помощью модели из раздела~\ref{sec:pqmethod}. Доля инфраструктуры
в приросте сообщения справедлива для одного уровня УЦ и для сертификатов аутентификации, выпущенных
тем же УЦ, что и сертификаты подписи.}

Таким образом, для RQ3 разделение переносится на постквантовые подписи\newred{, в виде оценки}. Их вычисления на сообщение
меняются примерно на миллисекунду, что меньше разброса между повторениями и меньше
архитектурных расходов в локальной сети. Их размер стоит круговых обходов, но только на новом
соединении\newred{: при испытанной конфигурации TCP} при добавленных 50\,мс один или два дополнительных
круговых обхода~--- это 50--100\,мс на сообщение против 153\,мс, которые повторное использование соединений уже экономит при
10\,КБ\newred{. При повторно используемых соединениях дополненным сообщениям не потребовался дополнительный круговой обход, и при
добавленных 50\,мс их медианная задержка составила от 139 до 145\,мс против 148\,мс для сообщений нынешнего
размера}.\par}

\section{Обсуждение}

\subsection{\newtext{Криптография и остаток тракта}}\label{sec:remainder}

\newtext{Таблица~\ref{tab:division} делит тракт X-Road на расходы, которые конфигурация может
изменить, и остаток в 22--27\,мс при 10\,КБ. Бенчмарк из RQ3 показывает, какую часть этого
остатка могут составлять подписи: подпись и проверка одного обмена заняли 1.7\,мс на ноутбуке
разработчика, где остаток составил 22.2\,мс, то есть криптографические операции~--- менее десятой
его части; в профиле постквантового TLS с взаимной аутентификацией криптография также составляла малую
часть, около 3\,\% процессорного времени~\cite{rahul2026}. Остальное~--- работа сервера безопасности, которую не отключает никакая конфигурация: журналирование записей
обмена без тел (19.3\,КБ на потребителе), формирование, разбор и каноникализация
сообщений и их подписей и их пересылка. Эти части нельзя разделить
методами этой работы, но вместе они составляют наибольшую часть тракта в локальной сети.
Поэтому разделение, о котором спрашивает RQ1, получается не двойным, а тройным: криптографические операции,
которые малы; архитектурные решения, которые можно изменить конфигурацией,~--- они малы в
локальной сети, но преобладают при больших телах и с расстоянием; и \newred{остаточная обработка} на
сервере безопасности вокруг криптографии, которая составляет наибольшую часть задержки небольшого сообщения
в локальной сети.} \newtext{Рис.~\ref{fig:summary} подытоживает, на что уходит время: в
локальной сети (a) и при добавленном времени кругового обхода 50\,мс для сообщений размера
постквантовых сообщений (b).}

\begin{figure*}[!tb]
\centering
\begin{tikzpicture}
\begin{axis}[name=pa, xbar stacked, width=0.44\textwidth, height=3.9cm, bar width=10pt,
  xmin=0, xmax=52, ytick=data, symbolic y coords={ответ 1\,МБ,ответ 10\,КБ},
  enlarge y limits=0.45, xlabel={Задержка тракта X-Road (мс)},
  title={(a) Локальная сеть}, title style={font=\footnotesize},
  tick label style={font=\scriptsize}, label style={font=\footnotesize},
  legend style={font=\scriptsize, at={(0.5,-0.42)}, anchor=north, legend columns=2,
    /tikz/every even column/.append style={column sep=0.6em}}, legend cell align=left]
\addplot[fill=okVermillion, draw=black] coordinates {(1.7,ответ 10\,КБ) (1.7,ответ 1\,МБ)};
\addlegendentry{подпись и проверка}
\addplot[fill=okSky!60, draw=black] coordinates {(20.5,ответ 10\,КБ) (27.4,ответ 1\,МБ)};
\addlegendentry{\newred{остаточная обработка}}
\addplot[fill=okOrange, draw=black] coordinates {(1.6,ответ 10\,КБ) (2.9,ответ 1\,МБ)};
\addlegendentry{новое соединение}
\addplot[fill=okGreen!70, draw=black] coordinates {(0.9,ответ 10\,КБ) (17.0,ответ 1\,МБ)};
\addlegendentry{журналирование тела}
\end{axis}
\begin{axis}[at={(pa.east)}, anchor=west, xshift=2.9cm, xbar, width=0.44\textwidth, height=5.2cm,
  bar width=5pt, xmin=0, xmax=650, ytick=data,
  symbolic y coords={C SLH-DSA-256s,C SLH-DSA-128f,C SLH-DSA-128s,B ML-DSA-87,RSA-2048 сейчас},
  enlarge y limits=0.13, xlabel={Медианная задержка при добавленных 50\,мс (мс)},
  title={(b) Добавленное время кругового обхода, дополненные сообщения}, title style={font=\footnotesize, align=center, text width=6cm},
  tick label style={font=\scriptsize}, label style={font=\footnotesize},
  legend style={font=\scriptsize, at={(0.5,-0.3)}, anchor=north, legend columns=2,
    /tikz/every even column/.append style={column sep=0.6em}}, legend cell align=left,
  nodes near coords, nodes near coords style={font=\tiny}, point meta=x]
\addplot[fill=okOrange, draw=black] coordinates {(309.4,RSA-2048 сейчас) (393.1,B ML-DSA-87) (393.6,C SLH-DSA-128s)
  (496.3,C SLH-DSA-128f) (541.8,C SLH-DSA-256s)};
\addlegendentry{новое соединение (\emph{full})}
\addplot[fill=okBlue!70, draw=black] coordinates {(147.8,RSA-2048 сейчас) (138.6,B ML-DSA-87) (138.8,C SLH-DSA-128s)
  (140.8,C SLH-DSA-128f) (145.0,C SLH-DSA-256s)};
\addlegendentry{повторное соединение (\emph{reuse})}
\end{axis}
\end{tikzpicture}
\caption{\newtext{На что уходит время обмена, ноутбук разработчика, один клиент. (a) Тракт X-Road
(\emph{full} $-$ \emph{direct}, таблица~\ref{tab:division}), разделённый на подпись и проверку
(бенчмарк, раздел~\ref{sec:rq3}), \newred{остаточную обработку на сервере безопасности (остаток тракта минус времена бенчмарка, получено вычитанием)}, открытие нового соединения (\emph{full} $-$ \emph{reuse}) и журналирование тела
(\emph{full} $-$ \emph{nobody}); синхронная простановка меток времени, не входящая в конфигурацию по умолчанию, добавила бы ещё
62--68\,мс. (b) Медианная задержка при добавленном времени кругового обхода 50\,мс для сообщений, дополненных до размера
нынешних сообщений и постквантовых сообщений в сценариях B и C (одна сетка,
таблица~\ref{tab:pqcheck}); постквантовые подпись и проверка изменили бы её на 0.2--6.4\,мс.}}
\label{fig:summary}
\end{figure*}

\subsection{Что результаты означают для предлагаемых изменений}\label{sec:pq}

Предложения, рассмотренные во введении, затрагивают разные механизмы, и измерения ставят
их в разное положение.

\newtext{Постквантовые подписи меняют подпись и проверку, которые оплачиваются на каждом сообщении
и выполняются на процессоре серверов безопасности. В поэтапной разбивке эту работу нельзя отделить от
остальной работы тех же этапов, но измеренная отдельно (раздел~\ref{sec:rq3}) она меняет
вычисления одного обмена примерно на миллисекунду для схем на решётках, что меньше разброса
между повторениями и меньше архитектурных расходов в локальной сети; только составные схемы
с RSA-3072 и RSA-4096 и SLH-DSA в УЦ добавляют от 1.6 до 9.3\,мс\newtext{, а на процессоре столь же медленном,
как на трёх машинах,~--- примерно до 19\,мс, что сопоставимо с остатком тракта
(раздел~\ref{sec:rq3})}. Размер важнее, и он
определяется не столько участниками, сколько инфраструктурой: каждое сообщение несёт сертификат
подписи и ответы OCSP, поэтому смена УЦ и ответчика OCSP увеличивает каждое
сообщение, и на неё приходится около трёх четвертей прироста (таблица~\ref{tab:pqD}). Ступень на
начальном окне~--- та же, что известна по постквантовым рукопожатиям
TLS~\cite{sikeridis2020,westerbaan2021,chou2026}; в X-Road она оплачивается на каждом сообщении на
новом соединении, а не один раз на рукопожатие. При политике соединений по умолчанию и расстоянии
между серверами безопасности сообщениям такого размера требуется до двух дополнительных круговых обходов на каждое с
ML-DSA, один или два, как только мигрируют УЦ и ответчик OCSP, и до пяти с SLH-DSA в УЦ, а при повторно используемых соединениях~--- ни одного. Размещение SLH-DSA в
УЦ, а не в ключах участников согласуется с размещением, признанным практичным для иерархий
TLS~\cite{delgado2026}.} Рукопожатие TLS между серверами
безопасности, которое постквантовые алгоритмы также увеличили бы~\cite{bakhtina2026pqc}, в каждом захваченном сообщении возобновлялось по
сеансовому билету без обмена сертификатами, так что более крупные сертификаты
затронули бы главным образом новые сеансы, а не каждое сообщение. \newtext{Между серверами безопасности
X-Road~7.8.3 использует только TLS~1.2~\cite{xroadsrc}, тогда как постквантовые обмен ключами и аутентификация
изучаются и стандартизируются для TLS~1.3~\cite{paquin2020,sikeridis2020}, поэтому постквантовый канал
сначала потребовал бы перехода X-Road на TLS~1.3. Полные рукопожатия тоже не редкость: потребитель, который
обменивается сообщениями со многими поставщиками или у которого истёк сеансовый билет, устанавливает новые сеансы,
и с постквантовыми сертификатами каждый из них нёс бы более крупные цепочки сертификатов, которые могут добавить
собственные круговые обходы~\cite{sikeridis2020}. Эта стоимость здесь не оценивается.}

Децентрализованная PKI меняет способ проверки сертификатов~\cite{bakhtina2023ares}. В нынешней
системе проверка статуса не входила в тракт запроса: сервер безопасности брал статус из
своего кэша, а запросы к ответчику OCSP появлялись, только когда на прогон приходилось фоновое
обновление. Поэтому любое предъявление
учётных данных на каждый запрос было бы новой стоимостью в тракте запроса, а не
заменой существующей.

Пороговые подписи меняют подпись~\cite{bakhtina2026ijis}. Сообщаемые для них круговые обходы от 191 до
272\,мс для небольших пороговых групп против 32\,мс с одним программным токеном в несколько раз превышают
всю стоимость X-Road, измеренную здесь с одним клиентом в локальной сети, и имеют тот же порядок, что
эффект политики соединений при добавленном времени кругового обхода 50\,мс.

\begin{table*}[!tb]
\centering\small
\caption{Предшествующие работы в сравнении с результатами этой работы.}
\label{tab:related}
\color{blue}
\begin{tabular}{@{}p{0.15\textwidth}p{0.38\textwidth}p{0.41\textwidth}@{}}
\toprule
Работа & Условия и результат & Эта работа \\
\midrule
Ansper et al.~\cite{ansper2013} & Проект пакетной простановки меток времени, аналитическая оценка;
реализация не измерялась & Около двух запросов к службе меток времени на 1200 сообщений;
синхронная политика добавляет 89--92\,мс и снижает потолок пропускной способности до 37--47\,\% \\
\addlinespace
Taklai~\cite{taklai2020} & Сервер безопасности перед сервисом реестра, 10--30 клиентов; преобладает время
базы данных; тракт через сервер безопасности быстрее прямого при 20 и 30
клиентах & Разброс под нагрузкой 15--41\,\% на трёх машинах, слишком велик для сравнения конфигураций;
приводятся только потолки и \emph{synctsa} \\
\addlinespace
Bakhtina et al.~\cite{bakhtina2026ijis} & X-Road 7.7.0, все компоненты на одном ноутбуке, среднее по 1000
запросов: 32\,мс с программным токеном, 191--272\,мс с пороговыми подписями & 25--30\,мс с одним
клиентом в трёх окружениях; один узел представителен для задержки в локальной сети, но не для
пропускной способности (в 1.4--1.8 раза выше на трёх машинах) и не для расстояния \\
\addlinespace
Snetkov and Vakarjuk~\cite{snetkov2023} & Постквантовые подписи в UXP: подпись Dilithium5
занимает 4595 байт против 256 у RSA-2048; времена не измерялись & Обмен передаёт около 40\,КБ, две такие
подписи добавили бы около 8.7\,КБ; сеансы TLS возобновляются без сертификатов; \newtext{с
сертификатами и ответами OCSP в каждом сообщении ML-DSA добавляет 4.6--41\,КБ и 0--2 круговых обхода
на новом соединении (таблица~\ref{tab:pq})} \\
\addlinespace
\newtext{Müller and Oupický~\cite{mueller2024}} & \newtext{Постквантовые и гибридные подписи XML в
Apache Santuario и BouncyCastle; время подписи и проверки и размер документа, без
цепочек сертификатов и сети} & \newtext{Те же библиотеки; на один обмен постквантовые
подписи меняют вычисления примерно на 1\,мс, тогда как их размер стоит круговых обходов на новых
соединениях} \\
\addlinespace
\newtext{Sikeridis et al.~\cite{sikeridis2020}, Chou and Cao~\cite{chou2026}} & \newtext{Постквантовые
цепочки сертификатов в TLS добавляют круговые обходы, как только превышают начальное окно TCP} &
\newtext{Та же ступень для каждого сообщения X-Road на новом соединении, а с четвёртой серии~---
ограничение окном приёма поставщика} \\
\addlinespace
\newtext{Vakarjuk et al.~\cite{vakarjuk2024}} & \newtext{Препятствия постквантовой миграции в
системах электронного правительства Эстонии; подписи для X-Road должны быть короткими; измерений нет} &
\newtext{Подписи инфраструктуры составляют около трёх четвертей прироста сообщения;
размер стоит круговых обходов только на новых соединениях} \\
\addlinespace
\newtext{Abbasi et al.~\cite{abbasi2025}, Rahul et al.~\cite{rahul2026}, Demir et
al.~\cite{demir2026}} & \newtext{Постквантовый (m)TLS~1.3: операции менее 1\,мс, менее
5\,\% добавленной задержки на серверах, около 3\,\% процессорного времени на криптографию; стоимость заключается в
размере рукопожатия} & \newtext{Тот же порядок в X-Road: подписи одного обмена занимают менее
десятой части остатка тракта, а размер увеличивает каждое сообщение, а не только
рукопожатие} \\
\addlinespace
Mironova~\cite{mironova2017} & Эстонский экземпляр в 2017 году: глобальная конфигурация действительна шесть часов,
свежесть OCSP~--- восемь часов & На стенде интервалы короче; фоновые расходы на сообщение в
промышленной эксплуатации были бы ниже \\
\addlinespace
Naylor et al.~\cite{naylor2014} & Стоимость TLS зависит главным образом от того, используются ли соединения повторно &
При добавленных 50\,мс повторное использование соединений экономит 54\,\% при 10\,КБ и 72\,\% при 1\,МБ \\
\bottomrule
\end{tabular}
\end{table*}

\subsection{Политика соединений}

Среди сравниваемых конфигураций наибольший эффект даёт решение открывать новое соединение для
каждого сообщения. В локальной сети оно стоит около 2\,мс и 6.2\,КБ на обмен (3.3\,КБ в сторону поставщика
и 2.9\,КБ обратно),
поэтому его легко не заметить. При добавленном времени кругового обхода каждое сообщение ожидает нескольких
дополнительных круговых обходов, и при 50\,мс повторное использование соединений экономит 54\,\% задержки для ответа
10\,КБ и 72\,\% для ответа 1\,МБ и почти удваивает пропускную способность с 16 клиентами. Это
согласуется с наблюдением, что стоимость TLS зависит главным образом от того, используются ли соединения
повторно~\cite{naylor2014}. Повторное использование соединений по умолчанию выключено~\cite{xroadsyspar}, так что его польза растёт
с расстоянием между серверами безопасности взаимодействующих организаций. Оно не лишено и
побочных эффектов: в измерениях единичные запросы по повторно используемым соединениям завершались ошибкой, что согласуется
с тем, что соединение из пула уже было закрыто другой стороной. \newtext{Документация
X-Road приводит и причину такого значения по умолчанию: без соединений из пула каждое соединение закрывается
после одного сообщения, что рекомендуется для серверов безопасности за балансировщиком нагрузки, чтобы он мог
равномерно распределять трафик~\cite{xroadsyspar,xroadxlb}. Возникали бы отказы, наблюдавшиеся с повторно используемыми
соединениями, и в промышленной эксплуатации и приемлема ли их частота, \newgreen{до} двух запросов на 10\,000
под нагрузкой, для сервисов, использующих X-Road, не исследовалось. Поэтому включение повторного использования~---
это компромисс для операторов, а не бесплатное улучшение.}

Политика соединений определяет и то, сколько стоят более крупные постквантовые сообщения. \newtext{При повторно используемых
соединениях окно TCP уже открыто, и ни одной из схем таблицы~\ref{tab:pq} не потребовался
дополнительный круговой обход. При новом соединении для каждого сообщения им требуется до двух дополнительных круговых
обходов с ML-DSA, Falcon и составными схемами и до пяти с SLH-DSA в УЦ; при добавленных
50\,мс один или два круговых обхода~--- это 50--100\,мс на сообщение против 153\,мс, которые повторное использование
соединений уже экономит для ответа 10\,КБ (\newgreen{раздел~\ref{sec:rtt}}).}

\subsection{Сравнение с предшествующими измерениями}

{\color{blue}В таблице~\ref{tab:related} результаты сопоставлены с предшествующими работами по X-Road и по стоимости TLS.\par}


Сквозная задержка с одним клиентом, от 25 до 30\,мс в трёх использованных здесь окружениях, имеет тот же
порядок, что 32\,мс, приведённые для X-Road~7.7.0 со всеми компонентами на одном ноутбуке~\cite{bakhtina2026ijis}.
То значение было представлено как нижняя граница, поскольку между серверами не было
сети~\cite{bakhtina2026ijis}. Для одного клиента в локальной сети сравнение одной машины с
тремя показывает, что такое значение, полученное на одном узле, представительно, поскольку стоимость сети
компенсируется тем, что у поставщика собственный процессор. Оно не представительно для пропускной способности, которая была
в 1.4--1.8 раза выше на трёх машинах, и для серверов безопасности, расположенных дальше друг от друга, где
преобладает политика соединений. В более раннем измерении сервера безопасности перед сервисом
реестра сообщалось, что тракт через сервер безопасности быстрее прямого при 20 и 30
одновременных клиентах~\cite{taklai2020}, что наводит на мысль, что изменчивость окружения была
сопоставима с измеряемым эффектом; и здесь разброс под нагрузкой был слишком велик для сравнения
конфигураций, тогда как с одним
клиентом и при добавленном времени кругового обхода он в основном оставался в пределах нескольких процентов.

Необходимый прогрев оказался гораздо длиннее, чем следует из предшествующих работ. В ближайшем измерении приводятся
средние по 1000 запросов без описания прогрева~\cite{bakhtina2026ijis}, тогда как здесь задержка не установилась
в пределах первой тысячи запросов, а первый прогон после перезапуска оставался более медленным даже после 4000, что согласуется с
рекомендацией измерять системы на Java только в установившемся режиме~\cite{georges2007}. Поэтому короткие измерения
X-Road могут завышать его задержку и сильнее различаться между прогонами.

\subsection{Абсолютные значения и разности}

Абсолютные значения зависят от оборудования: одна и та же сетка дала 24.8\,мс на ноутбуке разработчика и
28.3\,мс на трёх машинах для ответа 10\,КБ, а расходы, зависящие от процессора,~--- журналирование большого тела
и синхронная простановка меток времени~--- различались до двух раз. Разности между
конфигурациями сохраняли знак и порядок во всех окружениях, а эффект политики
соединений в локальной сети был примерно одинаковым (1.6--2.0\,мс при 10\,КБ). Режим питания процессора оказывал
ещё большее влияние, чем выбор машины. Поэтому при измерениях X-Road на оборудовании общего назначения
следует вместе с результатами указывать процессор и его режим питания.

\subsection{\newtext{Угрозы достоверности}}\label{sec:threats}

Несколько свойств стенда сдвигают абсолютные значения. Машины~--- ноутбуки, а не
серверы, а ноутбук разработчика добавляет накладные расходы WSL~2 в Windows. Служба меток времени
тестового УЦ заменена, а образы тестового УЦ и генератора нагрузки не были закреплены
за конкретной версией. Тестовый УЦ, ответчик OCSP
и служба меток времени работают рядом с серверами безопасности, тогда как в промышленной эксплуатации это
внешние службы с собственной задержкой, поэтому значения для проверки статуса и синхронной
простановки меток времени оптимистичны. Все интервалы фоновых механизмов оставлены по умолчанию,
тогда как промышленное развёртывание может использовать более длинные: для эстонского экземпляра 2017 года в магистерской
диссертации задокументированы срок действия глобальной конфигурации в шесть часов и свежесть OCSP в восемь
часов~\cite{mironova2017}, так что фоновые расходы на сообщение там были бы ещё ниже, чем на
стенде.

Другие свойства меняют не абсолютные значения, а доли и разброс. У эхо-сервиса
нет собственного времени обработки, что повышает относительную долю механизмов доверия по сравнению
с реальным сервисом. Под нагрузкой разброс между повторениями больше, поэтому там стабильная
пропускная способность и результаты синхронной простановки меток времени приводятся вместе с разбросом и
ошибками, а небольшие различия между остальными конфигурациями не приписываются отдельным
механизмам. Добавленное время кругового обхода эмулируется:
в нём нет постороннего трафика и потерь, и оно применяется внутри узлов серверов
безопасности\newtext{, что проверяется его применением к входящим, а не к исходящим пакетам
(раздел~\ref{sec:delay})}\newtext{; наклонов это не изменило (раздел~\ref{sec:rtt})}.
\newtext{Тела ответов~--- случайные байты, чтобы журнал сообщений не мог их сжать; реальные
сообщения JSON или XML сжимаются хорошо, поэтому стоимость журналирования большого тела~--- это наихудший случай. Ключи
участников находятся в программных токенах, тогда как промышленные серверы безопасности часто хранят их в аппаратных
модулях безопасности, где подпись идёт через PKCS\#11 и может занимать существенно больше времени. Причина
большого разброса пропускной способности под нагрузкой, например троттлинг процессоров ноутбуков, сборка мусора
виртуальной машины Java или база данных, не исследовалась.} \newred{Размеры 10\,КБ и
1\,МБ выбраны как небольшое и большое сообщение, а не по распределению размеров сообщений в реальном
экземпляре X-Road, которое не было доступно, и рабочая нагрузка использует один шаблон запросов к
эхо-сервису. Там, где конфигурация не дала разницы, это означает разницу меньше разброса между
повторениями, а для этапов~--- меньше разрешения 1\,мс, но не нулевую стоимость. Модель
серий справедлива для испытанной конфигурации TCP: другие начальные окна, алгоритмы управления перегрузкой,
потери пакетов или ограниченная полоса пропускания не исследовались.}

\newtext{Оценка для постквантовых подписей не является измерением X-Road с ними. Она
предполагает, что вычисления и круговые обходы складываются, что не проверялось.} \newtext{Дополненные
сообщения воспроизводят размер постквантовых сообщений, но не то, где находятся дополнительные байты: в
реальном сообщении они были бы в подписи XML, сертификатах и ответах OCSP, а не
в теле, и порядок, в котором X-Road записывает части сообщения в записи TLS, мог бы сдвинуть
границы серий. Не включают они и обработку более крупных структур внутри
X-Road, такую как разбор более крупных сертификатов и ответов OCSP, каноникализация XML с более длинными
значениями base64 и журналирование более крупных подписей, а также то, как постквантовые подписи в итоге были бы
встроены в X-Road, например как составная подпись или как две отдельные подписи, что могло бы изменить
и размер, и число операций.} \newtext{Оценённые} \newtext{круговые обходы} \newtext{зависят от настроек TCP на стенде, начального окна в 10 сегментов и
окна приёма поставщика, которое операционная система подстраивает сама, так что числа серий
справедливы для этой конфигурации; уменьшение сообщения, как с ECDSA в сценарии B, не проверялось, поскольку
дополнение может его только увеличить. \newviolet{Модель исправлялась дважды там, где она расходилась с
сетками, каждый раз по захватам трафика, и после последнего исправления сравнивалась только с
уже измеренными точками, а не с новыми. Поэтому для сообщения в пределах примерно одного сегмента
от границы серии оценка может ошибаться на один круговой обход.} \newgreen{Сетки проверки
выполнялись на ноутбуке разработчика, задержка которого без добавленной задержки менялась между сетками и внутри
них; поэтому проверка опирается на задержку при добавленном времени кругового обхода в сравнении с
базовым обменом той же сетки.} Предполагается, что сертификаты аутентификации используют схему
участников и тот же УЦ, а сертификаты рукопожатия TLS не учитываются, поскольку при возобновлённом
сеансе TLS ими не обмениваются. Вычисления измерены с
реализациями BouncyCastle~1.84 на Java~21, то есть версиями из X-Road~7.8.3. Собственная
ML-DSA есть в Java только начиная с Java~24~\cite{jep497}, и реализация из другой библиотеки или более поздней версии
Java может занимать иное время~\cite{oleksiichuk2026}.
Подписывающий компонент X-Road получает только хеш, а ML-DSA и SLH-DSA в чистом виде подписывают сами
данные, поэтому X-Road потребовались бы их варианты с предварительным хешированием~\cite{fips204,fips205} либо изменение этого
интерфейса; \newviolet{варианты с предварительным хешированием измерялись бенчмарком только для ML-DSA, SLH-DSA-128f и трёх
составных схем, но не для Falcon и не внутри подписывающего компонента X-Road}.
Составные схемы BouncyCastle~1.84 следуют черновикам, предшествовавшим RFC~10073; их идентификаторы объектов
совпадают с ним, но их подписи не сравнивались с его тестовыми векторами.}

Из-за этих факторов абсолютные задержки описывают этот стенд, а не промышленное
развёртывание. Поэтому сравнения проводятся между конфигурациями на одних и тех же машинах и
в одних и тех же условиях, а результаты следует понимать прежде всего как разности между конфигурациями
и как их изменение с размером сообщения, нагрузкой и добавленным временем кругового обхода; абсолютные значения приводятся
только для справки.

\section{Заключение}

В этой работе измерено, сколько стоят механизмы доверия X-Road в расчёте на одно сообщение: на трёх ноутбуках,
соединённых коммутатором, на одном из них, на более быстром ноутбуке и при управляемом времени кругового обхода,
добавленном между серверами безопасности.

По RQ1: тракт X-Road добавлял около 28\,мс к обмену размером 10\,КБ с одним клиентом на этих машинах.
Журналирование тела сообщения не давало измеримой разницы для небольших сообщений и составляло почти
половину задержки обмена размером 1\,МБ. Открытие нового соединения для каждого сообщения стоило около 2\,мс
и 6.2\,КБ на обмен в локальной сети. Проверка статуса сертификатов и пакетная простановка меток времени
не входили в тракт запроса \newred{в установившемся режиме, при действительных кэшированных ответах OCSP}: запросы к ответчику OCSP возникали только при фоновых обновлениях, а запросов меток времени было около двух
на 1200 сообщений,
тогда как синхронная простановка меток времени добавляла около 90\,мс. Подпись и проверку отключить нельзя;
\newtext{измеренные отдельно, они занимают 1.7\,мс на обмен на ноутбуке разработчика, менее
десятой части тех 22\,мс, которые остаются там от тракта после вычета расходов, изменяемых
конфигурацией. Поэтому разделение получается тройным: криптографические операции малы,
архитектурные решения, которые можно изменить конфигурацией, малы в локальной сети, но преобладают
при больших телах и с расстоянием, а \newred{остаточная обработка} на сервере безопасности вокруг
криптографии составляет наибольшую часть задержки небольшого сообщения в локальной сети.}

По RQ2: б\'{о}льшая часть роста при переходе от ответа 10\,КБ к ответу 1\,МБ приходилась на запись тела в
журнал сообщений. Под нагрузкой серверы безопасности достигали около 100 запросов в секунду на трёх машинах,
синхронная простановка меток времени \newtext{снижала это значение до 37--47\,\%}, а разброс между повторениями был слишком велик для сравнения
остальных конфигураций\newred{; поэтому для нагрузки работа отвечает на RQ2 только наибольшей
пропускной способностью, наблюдавшейся при 4, 16 и 32 клиентах, и эффектом синхронной простановки меток времени, а не
полными кривыми «пропускная способность~--- задержка»}.

\newtext{Также по RQ2}: одна машина и три машины дали почти одинаковые задержку и разности с одним
клиентом, тогда как три машины достигли большей пропускной способности. При добавленном времени кругового обхода
политика соединений по умолчанию заставляла каждое сообщение ожидать около пяти круговых обходов при 10\,КБ и десяти при 1\,МБ
против двух при повторно используемых соединениях; предсказание по одним лишь рукопожатиям давало три и один и
упускало два ожидания внутри TCP.

\newtext{По RQ3: постквантовые подписи, оценённые по бенчмарку на версии Java и
библиотеке X-Road и по модели серий сегментов TCP, проверенной на стенде, меняют вычисления
одного обмена примерно на миллисекунду для схем на решётках, что меньше архитектурных расходов,
и не более чем на 9\,мс для составных схем с RSA и при SLH-DSA в УЦ\newtext{, что на
более медленном процессоре трёх машин составило бы примерно до 19\,мс и сопоставимо с остатком
тракта}. Сообщениям их размера
не требуется дополнительный круговой обход при повторно используемых соединениях, но на новом соединении требуется один или два с ML-DSA,
как только мигрируют УЦ и ответчик OCSP, и до пяти с SLH-DSA в УЦ; б\'{о}льшая часть этого прироста приходится на
подписи инфраструктуры, которые несёт каждое сообщение.} \newviolet{Вычисления те же
и для вариантов с предварительным хешированием, которые потребовались бы подписывающему компоненту X-Road, получающему только
хеш.} \newred{Эти значения относятся к
подписям сообщений; постквантовые обмен ключами и сертификаты в рукопожатии TLS между
серверами безопасности, для которых X-Road сначала пришлось бы перейти на TLS~1.3, в оценку
не входят.}

Результаты дают точку отсчёта для нынешней системы, с которой можно сравнивать постквантовые подписи,
децентрализованную PKI и пороговые подписи. \newred{Они также указывают на то, что в
рассмотренных условиях и при испытанной конфигурации TCP стоимость более крупных постквантовых
подписей с точки зрения задержки определяется установлением соединения и окном TCP, а не криптографией: при
повторно используемых соединениях дополненным сообщениям не потребовался дополнительный круговой обход, тогда как при новом соединении
для каждого сообщения и далеко разнесённых серверах безопасности они стоят круговых обходов на каждом сообщении.} В дальнейшем можно измерить постквантовые
подписи, когда они станут доступны в X-Road, повторить измерения с внешними службами сертификации и
меток времени и по реальному каналу глобальной сети, а также использовать серверное оборудование.

\begin{thebibliography}{99}

\bibitem{mayer1995}
R.~C.~Mayer, J.~H.~Davis, F.~D.~Schoorman.
\newblock An Integrative Model of Organizational Trust.
\newblock \emph{Academy of Management Review}, 20(3):709--734, 1995.

\bibitem{grandison2000}
T.~Grandison, M.~Sloman.
\newblock A Survey of Trust in Internet Applications.
\newblock \emph{IEEE Communications Surveys \& Tutorials}, 3(4):2--16, 2000.

\bibitem{anderson2020}
R.~Anderson.
\newblock \emph{Security Engineering: A Guide to Building Dependable Distributed Systems}, 3rd ed.
\newblock Wiley, 2020.

\bibitem{bakhtina2023sac}
M.~Bakhtina, R.~Matulevi\v{c}ius, A.~Awad, P.~Kivim\"{a}ki.
\newblock On the Shift to Decentralised Identity Management in Distributed Data Exchange Systems.
\newblock In \emph{SAC '23}, pp.~864--873, ACM, 2023.
\newblock DOI: 10.1145/3555776.3577678.

\bibitem{xroadarcss}
Nordic Institute for Interoperability Solutions.
\newblock X-Road: Security Server Architecture.
\newblock \url{https://docs.x-road.global/Architecture/arc-ss_x-road_security_server_architecture.html}.
\newblock Дата обращения: 17 September 2026.

\bibitem{ansper2013}
A.~Ansper, A.~Buldas, M.~Freudenthal, J.~Willemson.
\newblock High-Performance Qualified Digital Signatures for X-Road.
\newblock In \emph{NordSec 2013}, LNCS 8208, pp.~123--138, Springer, 2013.

\bibitem{bakhtina2026ijis}
M.~Bakhtina, J.~Kvapil, P.~\v{S}venda, R.~Matulevi\v{c}ius.
\newblock Securing Organisational Digital Identity: A Full-Lifecycle Approach to Distributed Key
Management.
\newblock \emph{International Journal of Information Security}, 25:131, 2026.
\newblock DOI: 10.1007/s10207-026-01291-5.

\bibitem{taklai2020}
E.~Taklai.
\newblock X-tee teenuste j\~{o}udlustestimine Eesti \"{a}riregistri n\"{a}itel
(Performance testing of X-Road services on the example of the Estonian Business Register).
\newblock Diploma thesis, Tallinn University of Technology, 2020.

\bibitem{xteefactsheet}
Estonian Information System Authority (RIA).
\newblock X-tee Factsheet, Estonia.
\newblock \url{https://www.x-tee.ee/factsheets/EE/}.
\newblock Figures based on data of 1 September 2026; дата обращения: 17 September 2026.

\bibitem{bakhtina2026pqc}
M.~Bakhtina.
\newblock Migration of X-Road to Post-Quantum Cryptography.
\newblock NIIS Publication series, Nordic Institute for Interoperability Solutions, 2026.
\newblock ISBN 978-9916-9853-4-2.

\bibitem{snetkov2023}
N.~Snetkov, J.~Vakarjuk.
\newblock Integrating Post-Quantum Cryptography to UXP.
\newblock Cybernetica Research Report D-2-499, 2023.

\bibitem{vakarjuk2024}
\textcolor{blue}{J.~Vakarjuk, N.~Snetkov, P.~Laud.
\newblock Identifying Obstacles of PQC Migration in e-Estonia.
\newblock In \emph{16th International Conference on Cyber Conflict (CyCon 2024)}, pp.~63--81,
NATO CCDCOE, 2024.
\newblock DOI: 10.23919/CyCon62501.2024.10685570.}

\bibitem{bakhtina2023ares}
M.~Bakhtina, K.~L.~Leung, R.~Matulevi\v{c}ius, A.~Awad, P.~\v{S}venda.
\newblock A Decentralised Public Key Infrastructure for X-Road.
\newblock In \emph{ARES 2023}, ACM.
\newblock DOI: 10.1145/3600160.3605092.

\bibitem{coarfa2006}
\textcolor{blue}{C.~Coarfa, P.~Druschel, D.~S.~Wallach.
\newblock Performance Analysis of TLS Web Servers.
\newblock \emph{ACM Transactions on Computer Systems}, 24(1):39--69, 2006.
\newblock DOI: 10.1145/1124153.1124155.}

\bibitem{naylor2014}
D.~Naylor, A.~Finamore, I.~Leontiadis, Y.~Grunenberger, M.~Mellia, M.~Munaf\`{o},
K.~Papagiannaki, P.~Steenkiste.
\newblock The Cost of the ``S'' in HTTPS.
\newblock In \emph{CoNEXT '14}, pp.~133--140, ACM, 2014.
\newblock DOI: 10.1145/2674005.2674991.

\bibitem{paquin2020}
C.~Paquin, D.~Stebila, G.~Tamvada.
\newblock Benchmarking Post-quantum Cryptography in TLS.
\newblock In \emph{PQCrypto 2020}, LNCS vol.~12100, pp.~72--91, Springer, 2020.
\newblock Cryptology ePrint Archive 2019/1447.

\bibitem{sikeridis2020}
\textcolor{blue}{D.~Sikeridis, P.~Kampanakis, M.~Devetsikiotis.
\newblock Post-Quantum Authentication in TLS 1.3: A Performance Study.
\newblock In \emph{NDSS 2020}, Internet Society, 2020.
\newblock DOI: 10.14722/ndss.2020.24203.}

\bibitem{westerbaan2021}
\textcolor{blue}{B.~Westerbaan.
\newblock Sizing Up Post-Quantum Signatures.
\newblock Cloudflare blog, 8 November 2021.
\newblock \url{https://blog.cloudflare.com/sizing-up-post-quantum-signatures/}.}
\newblock \newtext{Дата обращения: 19 November 2026.}

\bibitem{chou2026}
\textcolor{blue}{M.~Chou, P.~Cao.
\newblock Network Impact of Post-Quantum Certificate Chain Sizes on Time to First Byte in TLS
Deployments.
\newblock arXiv:2604.24869, 2026.}

\bibitem{kampanakis2024}
\textcolor{blue}{P.~Kampanakis, W.~Childs-Klein.
\newblock The Impact of Data-Heavy, Post-Quantum TLS 1.3 on the Time-To-Last-Byte of Real-World
Connections.
\newblock In \emph{MADWeb 2024}, Internet Society, 2024.
\newblock DOI: 10.14722/madweb.2024.23010.}

\bibitem{abbasi2025}
\newtext{M.~Abbasi, F.~Cardoso, P.~V{\'a}z, J.~Silva, P.~Martins.
\newblock A Practical Performance Benchmark of Post-Quantum Cryptography Across Heterogeneous
Computing Environments.
\newblock \emph{Cryptography}, 9(2):32, 2025.
\newblock DOI: 10.3390/cryptography9020032.}

\bibitem{rahul2026}
\newtext{Rahul, M.~K.~Hanawal, S.~G.~Kulkarni, Saravanan~V, K.~T.~Reddy.
\newblock PQC Migration: End-to-End Performance Profiling of Post Quantum Cryptographic Algorithms in
mTLS.
\newblock arXiv:2611.21206, 2026.}

\bibitem{demir2026}
\newtext{E.~D.~Demir, I.~Gorgulu, M.~C.~Onba{\c{s}}l{\i}.
\newblock Micro and Transport Layer Security Benchmarking of Post-Quantum Versus Classical
Cryptography.
\newblock \emph{IEEE Access}, 14, 2026.
\newblock DOI: 10.1109/ACCESS.2026.3682760.}

\bibitem{mueller2024}
\textcolor{blue}{J.~M{\"u}ller, J.~Oupick{\'y}.
\newblock Post-quantum XML and SAML Single Sign-On.
\newblock \emph{Proceedings on Privacy Enhancing Technologies}, 2024(4):525--543, 2024.
\newblock DOI: 10.56553/popets-2024-0128.}

\bibitem{xroadsyspar}
Nordic Institute for Interoperability Solutions.
\newblock X-Road: System Parameters (UG-SYSPAR)\newtext{, version 2.103 of 24 March 2026}.
\newblock \url{https://docs.x-road.global/Manuals/ug-syspar_x-road_v6_system_parameters.html}.
\newblock Дата обращения: 17 September 2026.

\bibitem{kivimaki2020}
P.~Kivim\"{a}ki.
\newblock Resisting Failure.
\newblock NIIS blog, \newtext{26} November 2020.
\newblock \url{https://www.niis.org/blog/2020/11/20/resisting-failure}.
\newblock \newtext{Дата обращения: 19 November 2026.}

\bibitem{xroadsrc}
\newtext{Nordic Institute for Interoperability Solutions.
\newblock X-Road source code, version 7.8.3.
\newblock \url{https://github.com/nordic-institute/X-Road/tree/7.8.3}.
\newblock Дата обращения: 19 November 2026.}

\bibitem{k6}
\newtext{Grafana Labs.
\newblock Grafana k6, load testing tool.
\newblock \url{https://k6.io/}.
\newblock Дата обращения: 19 November 2026.}

\bibitem{chrony}
\newtext{The chrony project.
\newblock chrony, an implementation of the Network Time Protocol.
\newblock \url{https://chrony-project.org/}.
\newblock Дата обращения: 19 November 2026.}

\bibitem{hemminger2005}
S.~Hemminger.
\newblock Network Emulation with NetEm.
\newblock In \emph{linux.conf.au 2005}, Canberra, 2005.

\bibitem{nussbaum2009}
L.~Nussbaum, O.~Richard.
\newblock A Comparative Study of Network Link Emulators.
\newblock In \emph{Communications and Networking Simulation Symposium (CNS '09)}, 2009.

\bibitem{jurgelionis2011}
A.~Jurgelionis, J.-P.~Laulajainen, M.~Hirvonen, A.~I.~Wang.
\newblock An Empirical Study of NetEm Network Emulation Functionalities.
\newblock In \emph{ICCCN 2011}, IEEE, 2011.

\bibitem{georges2007}
A.~Georges, D.~Buytaert, L.~Eeckhout.
\newblock Statistically Rigorous Java Performance Evaluation.
\newblock In \emph{OOPSLA '07}, pp.~57--76, ACM, 2007.
\newblock DOI: 10.1145/1297027.1297033.

\bibitem{fips204}
National Institute of Standards and Technology.
\newblock Module-Lattice-Based Digital Signature Standard.
\newblock FIPS 204, 2024.
\newblock DOI: 10.6028/NIST.FIPS.204.

\bibitem{rfc10073}
M.~Ounsworth, J.~Gray, M.~Pala, J.~Klau{\ss}ner, S.~Fluhrer.
\newblock Composite Module-Lattice-Based Digital Signature Algorithm (ML-DSA) for Use in X.509
Public Key Infrastructure.
\newblock RFC 10073, IETF, 2026.

\bibitem{fips205}
National Institute of Standards and Technology.
\newblock Stateless Hash-Based Digital Signature Standard.
\newblock FIPS 205, 2024.
\newblock DOI: 10.6028/NIST.FIPS.205.

\bibitem{delgado2026}
\textcolor{blue}{J.~L.~Delgado Jim{\'e}nez.
\newblock Signature Placement in Post-Quantum TLS Certificate Hierarchies: An Experimental Study of
ML-DSA and SLH-DSA in TLS 1.3 Authentication.
\newblock arXiv:2604.06100, 2026.}

\bibitem{bouncycastle}
\newtext{Legion of the Bouncy Castle.
\newblock Bouncy Castle Crypto APIs for Java, version 1.84.
\newblock \url{https://www.bouncycastle.org/download/bouncy-castle-java/}.
\newblock Дата обращения: 19 November 2026.}

\bibitem{oleksiichuk2026}
\textcolor{blue}{Yu.~Oleksiichuk.
\newblock Empirical Performance Evaluation of NIST Post-Quantum Cryptographic Algorithms ML-KEM, ML-DSA,
and SLH-DSA across JDK Versions.
\newblock \emph{Cybersecurity: Education, Science, Technique}, 1(33), 2026.
\newblock DOI: 10.28925/2663-4023.2026.33.1210.}

\bibitem{jmh}
\newtext{OpenJDK.
\newblock Java Microbenchmark Harness (JMH).
\newblock \url{https://github.com/openjdk/jmh}.
\newblock Дата обращения: 19 November 2026.}

\bibitem{rfc9110}
\newtext{R.~Fielding, M.~Nottingham, J.~Reschke.
\newblock HTTP Semantics.
\newblock RFC 9110, IETF, 2022.}

\bibitem{mironova2017}
L.~Mironova.
\newblock X-tee platvormi kriitiliste kesksete teenuste testimine
(Testing of the critical central services of the X-Road platform).
\newblock Master's thesis, Tallinn University of Technology, 2017.

\bibitem{xroadxlb}
\newtext{Nordic Institute for Interoperability Solutions.
\newblock X-Road: External Load Balancer Installation Guide (IG-XLB), version 1.30.
\newblock \url{https://docs.x-road.global/Manuals/LoadBalancing/ig-xlb_x-road_external_load_balancer_installation_guide.html}.
\newblock Дата обращения: 19 November 2026.}

\bibitem{jep497}
\textcolor{blue}{OpenJDK.
\newblock JEP 497: Quantum-Resistant Module-Lattice-Based Digital Signature Algorithm.
\newblock Delivered in JDK 24.
\newblock \url{https://openjdk.org/jeps/497}.}
\newblock \newtext{Дата обращения: 19 November 2026.}

\end{thebibliography}

\end{document}
